%%%%%%% DO NOT REMOVE COMMANDS OR PACKAGE FROM THE DOCUMENT

\documentclass[letterpaper,10pt]{article}

% \usepackage[43-48]{pagesel}
% Disabled to support full-book and per-chapter compiles without page filtering.
%Restricts the pages being compiled

% week 1, 1-5

% week 2, 6
\usepackage{hyperref}
\usepackage{alltt}
\usepackage{listings}
\usepackage{amsfonts}
\usepackage{amsmath}
\allowdisplaybreaks[1]%0-4 4 is most permissive of breaking equations
\usepackage{amssymb}
\usepackage{amsthm}
\usepackage{array}
\usepackage{calc}
\usepackage{color}
\usepackage{graphicx}
\usepackage{stmaryrd}
\usepackage{supertabular}
\usepackage{url}%Wanted by the bibilography
\usepackage[all]{xy}

%%Packages for tables
\usepackage{wrapfig}
\usepackage{multirow}
\usepackage{tabu}
\usepackage{tabularx}

%Note: Throughout the preamble there are occassional references to pages out of
%books. Unless otherwise noted they are from The LaTeX Companion 2nd ed by
%Mittelbach and Goossens

%######################################
%Font Selection.
%Concrete
%\usepackage{beton}
%\usepackage{euler}
%Computer Modern Bright
%\usepackage[T1]{fontenc}
%\usepackage{cmbright}

%######################################

%######################################
%Quote environment
\newsavebox{\quauth}
\newenvironment{Quotation}[1]
	{\sbox{\quauth}{\textit{#1}}\begin{quote}}
	{\hspace*{\fill}\nolinebreak[1]\hspace*{\fill}\usebox{\quauth}\end{quote}}
%######################################

%######################################
%Theorem setup, needs package amsthm
%Theorems
\theoremstyle{plain}
\newtheorem{thm}{Theorem}[section]
\newtheorem{algo}[thm]{Algorithm}
\newtheorem{cor}[thm]{Corollary}
\newtheorem{lem}[thm]{Lemma}
\newtheorem{prop}[thm]{Proposition}


%Definitions
\theoremstyle{definition}
\newtheorem{appl}[thm]{Application}
\newtheorem{conj}[thm]{Conjecture}
\newtheorem{defn}[thm]{Definition}
\newtheorem{exmp}[thm]{Example}
\newtheorem{exer}[thm]{Exercise}

%Remarks
\theoremstyle{remark}
\newtheorem{rem}[thm]{Remark}
\newtheorem{note}[thm]{Note}
\newtheorem*{case}{Case}
\newtheorem*{claim}{Claim}


%For referencing Theorems/Defs/etc...
\providecommand{\thmref}[1]{ (Thm. \ref{#1})}
\providecommand{\defref}[1]{ (Def. \ref{#1})}
\providecommand{\exref}[1]{Example \ref{#1}}
%######################################

%######################################
%Define a few column types for convenience
%Needs package array

%These make the columns all be in math mode
%The extra %stopzone are to make vim do syntax
%highlighting correctly with this code, otherwise
%it chokes
\newcolumntype{C}{>{$}c<{$}}
%stopzone%stopzone%stopzone
\newcolumntype{L}{>{$}l<{$}}
%stopzone%stopzone%stopzone
\newcolumntype{R}{>{$}r<{$}}
%stopzone%stopzone%stopzone
%######################################

%######################################
%Various shortcuts
%Some may need package amssymb

%Shortcuts for various sets
\providecommand{\A}{\ensuremath{\mathfrak{A}}}%ideal A
\providecommand{\B}{\ensuremath{\mathcal{B}}}%basis
\providecommand{\C}{\ensuremath{\mathbb{C}}}%complex
\providecommand{\F}{\ensuremath{\mathbb{F}}}%finite field
\providecommand{\N}{\ensuremath{\mathbb{N}}}%natural
\providecommand{\RI}{\ensuremath{\mathcal{O}}}%ring of integers
\providecommand{\p}{\ensuremath{\mathfrak{P}}}%prime ideal
\providecommand{\Q}{\ensuremath{\mathbb{Q}}}%rationals
\providecommand{\R}{\ensuremath{\mathbb{R}}}%reals
\providecommand{\T}{\ensuremath{\mathcal{T}}}%topology
\providecommand{\Z}{\ensuremath{\mathbb{Z}}}%integers
\providecommand{\Zpos}{\ensuremath{\mathbb{Z}^{+}}}%positive integers


%Miscellaneous shortcuts
\providecommand{\lcm}{\ensuremath{\text{lcm}}}%least common multiple
\providecommand{\st}{\ensuremath{\backepsilon}}
\providecommand{\ud}{\ensuremath{\,\,\mathrm{d}}}%For derivatives
%######################################

%######################################
% Various operators
% some may need package amsmath

%Algebra
\providecommand{\amod}{\ensuremath{\diagup}}%Algebraic mod
\providecommand{\embed}{\ensuremath{\hookrightarrow}}%inclusion map
\providecommand{\gen}[1]{\ensuremath{\left<#1\right>}}%Group generated by #1
\providecommand{\idx}[2]{\ensuremath{\left[#1:#2\right]}}%index of #2 in #1
\providecommand{\im}{\ensuremath{\mathrm{im}}}%Image
\providecommand{\iso}{\ensuremath{\simeq}}%isomorphism
\providecommand{\normal}{\ensuremath{\vartriangleleft}}%Normal subgroup
\providecommand{\subgp}{\ensuremath{\le}}%Sub group
\providecommand{\vect}[1]{\ensuremath{\left\langle#1\right\rangle}}%Vector

%Logic
\providecommand{\land}{\ensuremath{\wedge}}
\DeclareMathSymbol\lnot{\mathbin}{symbols}{"3A}%p528
\providecommand{\lor}{\ensuremath{\vee}}
\providecommand{\lxor}{\ensuremath{\oplus}}
\providecommand{\todo}[1]{\textcolor{magenta}{\textbf{#1}}}

%Number Theory
\providecommand{\jacobi}[2]{\ensuremath{\left(\frac{#1}{#2}\right)}}
\providecommand{\legendre}[2]{\ensuremath{\left(\frac{#1}{#2}\right)}}

%Set Theory
\providecommand{\intersect}{\ensuremath{\bigcap}}
\providecommand{\less}{\ensuremath{\diagdown}}
\providecommand{\union}{\ensuremath{\bigcup}}

%Miscellaneous
\providecommand{\abs}[1]{\ensuremath{\left\lvert#1\right\rvert}}
\providecommand{\card}[1]{\ensuremath{\left\lvert#1\right\rvert}}
\providecommand{\ceiling}[1]{\ensuremath{\left\lceil#1\right\rceil}}
\providecommand{\define}{\ensuremath{\stackrel{\text{\tiny def}}{=}}}
\providecommand{\floor}[1]{\ensuremath{\left\lfloor#1\right\rfloor}}
\providecommand{\norm}[1]{\ensuremath{\left\lVert#1\right\rVert}}%Analysis norm
\providecommand{\order}[1]{\ensuremath{\left\lvert#1\right\rvert}}

%Asymptotic
\providecommand{\Oh}{\ensuremath{\mathcal{O}}}
\providecommand{\oh}{\ensuremath{o}}
%######################################
%######################################
%Page layout
% If using hyperef, load it before this block
\usepackage[paper=letterpaper,tmargin=1in,bmargin=42pt,lmargin=.5in,rmargin=.5in,headheight=30pt,headsep=30pt,footskip=20pt]{geometry}
% Note 1in = 72pt, therefore bmargin=42pt is 1in - 30pt
\usepackage{ifpdf}
\ifpdf
  \geometry{driver=pdftex}
\else
  \geometry{driver=dvips}
\fi

%%%%%%%%%%%%
%%%%%%%%%%%%
%%%%%%%%%%%%
%%%%%%%%%%%% Here you enter the running title
%
%%%%%%%%%%%%
\usepackage{fancyhdr}
\lhead{COMP 2300}
%It is recommended to add your name
%%%%%%%%%%%%
\chead{Notes}
\rhead{Fall 2025}
\lfoot{}
\cfoot{}
\rfoot{\thepage}
\renewcommand{\headrulewidth}{0pt}
\renewcommand{\footrulewidth}{0pt}
%######################################
%######################################
%######################################
%######################################
%######################################


% Compile a single chapter by defining:
% \def\coursehighlightsincludeonly{chapters/01_foundations}
\ifdefined\coursehighlightsincludeonly
  \includeonly{\coursehighlightsincludeonly}
\fi

\begin{document}
\pagestyle{fancy}

\section{The Foundations}
\subsection{Logic and Proofs}

\begin{defn}[proposition]
A \textbf{proposition} is a declarative sentence (that is, a sentence that declares a fact) that is either true
or false, but not both.
\end{defn}

\begin{defn}[propositional variables]
We use letters to denote \textbf{propositional variables} (or statement variables), that is, variables that represent propositions, just as letters are used to denote numerical variables.
\end{defn}
\begin{defn}[truth value] The \textbf{truth value} of a proposition is true, denoted by $T$, if it is a true proposition, and the truth value of a proposition
is false, denoted by $F$, if it is a false proposition.
\end{defn}

\begin{defn}[negation]
Let $p$ be a proposition. The \textbf{negation} of $p$, denoted by $ \lnot p$, is the statement 
\newline\todo{"It is not the case that p"}. 
\newline The proposition $\lnot p$ is read “not p.” The truth value of the negation of $p$, $\lnot p$, is the opposite of the truth value of $p$.
\end{defn}

\begin{exmp}[truth table of $\lnot$]
Below is an example of a truth table. 

Typically, the leftmost columns correspond to the propositional variables that we use in our model. 

Each row in the table corresponds to a different assignment of the propositional variables, each possible combination of T and F values.

For example, the first row below corresponds to the world (or model) in which $p$ is True. The second row corresponds to world (or model) in which $p$ is False. 

The columns to the right of the propositional variables tell us the value of $\lnot p$ corresponding to the truth value of $p$ in this row.

\begin{center}
\begin{tabular}{ | c | c | c |} 
  \hline
 $p$& $ \lnot p$ \\
 \hline
  T & F \\ 
  F & T \\ 
  \hline
\end{tabular}
\end{center}
\end{exmp}

Suppose that $p$ is the proposition "candy is free"


Then the first row corresponds to a world in which the proposition "candy is free" is True, in which case the negation of this proposition "it is not the case that candy is free" is False.

The second row corresponds to a world in which the proposition "candy is free" is \todo{?False}, in which case the negation of this proposition "it is not the case that candy is free" is \todo{True}

\begin{defn}[conjunction,($\land$)]
Let $p$ and $q$ be propositions. The \textbf{conjunction} of $p$ and $q$, denoted by $p \land q$, is the proposition
“$p$ and $q$.” The conjunction $p \land q$ is True when both $p$ and $q$ are True and is False otherwise.
\end{defn}

\begin{exmp}[truth table of $\land$]
Below is the \textbf{truth table} of $p \land q$. In the leftmost columns are the propositional values $p$ and $q$. Each row represents a different world (or model or assignment) in which $p$ and $q$ take on the values True and False. 

\begin{center}
\begin{tabular}{ | c | c | c |} 
  \hline
 $p$ $q$ & $ p \land q$ \\
 \hline
  T T & T \\ 
  T F & F \\ 
  F T & F \\ 
  F F & F \\ 
  \hline
\end{tabular}
\end{center}

The first row represent a world in which both $p$ and $q$ are true. 

The second row represents a world in which $p$ is True and $q$ is \todo{F}.

The third row represents a world in which $p$ is \todo{F} and $q$ is \todo{T}.

The last row represents \todo{a world in which p and q are both false}

Suppose that $p$ represents the proposition "candy is free" and that $q$ represents the proposition "we all get along".
Then the first row represents a world in which $p$ and $q$ are true, in which case the proposition "candy is free and we all get along " is true. The second row, however, represents a world in which the proposition "candy is free" is \todo{True} and "we all get along is" \todo{False} in which case the proposition "candy is free and we all get along " is \todo{False}.
\end{exmp}


\newpage
\begin{defn}[disjunction $\lor$]
Let $p$ and $q$ be propositions. The \textbf{disjunction} of $p$ and $q$, denoted by $p \lor q$, is the proposition
“p or q.” The disjunction $p \lor q$ is false when both p and q are false and is true otherwise.
\end{defn}

\begin{exmp}[truth table of or $\lor$]
Below is the truth table of $p \lor q$. In the leftmost columns are the propositional values $p$ and $q$. Each row represents a different world (or model or assignment) in which $p$ and $q$ take on the values True and False. 

\begin{center}
\begin{tabular}{ | c | c | c |} 
  \hline
 $p$ $q$ & $ p \lor q$ \\
 \hline
  T T & \todo{T} \\ 
  T F & \todo{T} \\ 
  F T & \todo{T} \\ 
  F F & \todo{F} \\ 
  \hline
\end{tabular}
\end{center}

Suppose that $p$ represents the proposition "I am going to win my soccer match" and that $q$ represents the proposition "I am going to cry".

The first row represent a world in which both $p$ and $q$ are true.
This means that I will \todo{win} my soccer match and then I will \todo{cry}. The proposition "I am going to win my soccer match or I am going to cry" is \todo{True}.

The second row represents a world in which $p$ is True and $q$ is False. This means that I will \todo{Win} my soccer match and then I will \todo{Not Cry}. The proposition "I am going to win my soccer match or I am going to cry" is \todo{True}.

The third row represents a world in which $p$ is \todo{False} and $q$ is \todo{True}. This means that I will \todo{Lose} my soccer match and then I will \todo{Cry}. The proposition "I am going to win my soccer match or I am going to cry" is \todo{True}.

The last row represents a world in which $p$ is \todo{False} and $q$ is \todo{False}. This means that I will \todo{Lose} my soccer match and then I will \todo{Not Cry}. The proposition "I am going to win my soccer match or I am going to cry" is \todo{False}.
\end{exmp}
\begin{defn}[exclusive or $\lxor$]
Let $p$ and $q$ be propositions. The \textbf{exclusive or} of $p$ and $q$, denoted by $p \lxor q$, is the proposition
that is true when exactly one of $p$ and $q$ is true and is false otherwise.
\end{defn}
\begin{exmp}[truth table of $\lxor$]

\begin{center}
\begin{tabular}{ | c | c | c |} 
  \hline
 $p$ $q$ & $ p \lxor q$ \\
 \hline
  T T & \todo{F} \\ 
  T F & \todo{T} \\ 
  F T & \todo{T} \\ 
  F F & \todo{F} \\ 
  \hline
\end{tabular}
\end{center}
\end{exmp}
If $p$ is the statement "My side is soup" and $q$ is the statement "My side is salad", then $p \lxor q$ represents the proposition "My side is soup or salad (but not both)", which is true in models where exactly one of the propositions "my side is soup" and "my side is salad" is true.

\begin{defn}[conditional $p \to q$]
Let $p$ and $q$ be propositions. The conditional statement $p \to q$ is the proposition “if $p$, then
$q$.” The conditional statement $p \to q$ is false when $p$ is true and $q$ is false, and true otherwise.
In the conditional statement $p \to q$, $p$ is called the hypothesis (or antecedent or premise)
and $q$ is called the conclusion (or consequence).
\end{defn}
\begin{exmp}[truth table of $\to$]
\begin{center}
\begin{tabular}{ | c | c | c |} 
  \hline
 $p$ $q$ & $ p \to q$ \\
 \hline
  T T & T \\ 
  T F & F \\ 
  F T & T \\ 
  F F & T \\ 
  \hline
\end{tabular}
\end{center}
If $p$ is the proposition "You give me 1000 dollars" and $q$ is the statement "Your roof will stop leaking." Then $p \to q$ represents the proposition, 
\newline 
\todo{"If you give me 1000 dollars your roof  will stop leaking"}.
\newline
The proposition $p\to q$ is False when the proposition "You give me 1000 dollars" is \todo{True} and "Your roof will stop leaking." is \todo{False}.
\end{exmp}

\newpage

\begin{exmp}[conditional $p\to q$ ] \hspace{\fill}

Consider the two propositions and their respective propositional variables

$ p:=\text{"You eat 20 hot dogs for dinner"} $

$ q:=\text{"You don't need a little snacky poo"} $

The proposition (or statement)
\newline
$p\to q:= $\todo{If you eat 20 hot dogs for dinner then you don't need a little snacky poo }
\newline
In a world in which you only ate 19 hotdogs for dinner and you don't want a snack, $p\to q$ is \todo{True}.

In a world in which you eat 20 hotdogs and you do still need a little snack, $p \to q$ is \todo{False}.

\end{exmp}

\begin{defn}[biconditional $p\iff q$]\ \\
Let $p$ and $q$ be propositions. The \textbf{biconditional} statement $p \iff q$ is the proposition “$p$ if
and only if $q$.” The biconditional statement $p \iff q$ is true when $p$ and $q$ have the same truth
values, and is false otherwise. Biconditional statements are also called bi-implications.
\end{defn}

\begin{exmp}[truth table of $\iff$]
\begin{center}
\begin{tabular}{ | c | c | c |} 
  \hline
 $p$ $q$ & $ p \iff q$ \\
 \hline
  T T & \todo{T} \\ 
  T F & \todo{F} \\ 
  F T & \todo{F} \\ 
  F F & \todo{T} \\ 
  \hline
\end{tabular}
\end{center}
If $p$ is the proposition "Harry potter dies." and $q$ is the statement "Voldemort dies." Then $p \iff q$ represents the proposition, \todo{"Harry potter dies if and only if Voldemort dies"}.

The proposition $p\iff q$ is False when the proposition "Harry potter dies." is True and "Voldemort dies." is \todo{False} OR when the proposition "Harry potter dies." is False and "Voldemort dies." is \todo{True}.

\end{exmp}

\newpage
\subsection{Applications of Propositional Logic}

\subsubsection{Logic Puzzles}
\begin{exmp}[logic puzzle]

On an island in the future there are two groups of people. One group always lies and the other group always tells the truth.

You come across 2 people.
One person says "I don't lie, but she does!"
The other person says "I don't lie, but she does!"

What are possible scenarios of who is the liar and who tells the truth?

Can you write the scenario above as a compound proposition?

\begin{proof} \ \\
I will interpret "I don't lie, but she does!" as equivalent to the statement "I don't lie \textbf{and} she does!"

$A:=\text{person 1 tells the truth}$

$B:=\text{person 2 tells the truth}$

Then person one says, $A\land (\lnot B)$
and person two says $B\land (\lnot A)$

Supposing that $A$ is true (person 1 is telling the truth), $A\land (\lnot B)$ must be true, in which case $\lnot B $ is true. This means person 1 is telling the truth and person 2 is lying. 
We now need to confirm that person 2 lied. If person 2 lied, $B$ is false and $B \land \lnot(A)$ must also be false. Therefore, $\lnot(B \land \lnot(A))$ must be true. Using De Morgans law, $\lnot B \lor \lnot A$ must be True, and since $\lnot B$ we see that this is the case. Therefore it is possible that person 1 told the truth and person 2 lied.

Supposing that $A$ is false (person 1 is lying), $A\land (\lnot B)$ must be false, which is satisfied by $A$ being false. Therefore $B$ may take the value false and the statement $A\land (\lnot B)$ is false (because person lied in the first half of the sentence and the best interpretation of 'but' in this context is 'and' but it makes me want to cry. Not all sentences translate perfectly.)

\end{proof}

\end{exmp}

\newpage

\subsection{Propositional Equivalences}

\begin{defn}[tautology and contradiction]

A compound proposition that is always true, no matter what the truth values of the propositional
variables that occur in it, is called a \textbf{tautology}. A compound proposition that is always
false is called a \textbf{contradiction}. A compound proposition that is neither a tautology nor a
contradiction is called a contingency.

\end{defn}
\begin{defn}[equivalence]
The compound propositions $p$ and $q$ are called logically equivalent if $p \iff q$ is a tautology.
The notation $p \equiv q$ denotes that p and q are logically equivalent. This means that the two statements evaluate identically in every possible world!
\end{defn}
\begin{exmp}[tautology]
A tautology is logically equivalent to True, while a contradiction is logically equivalent to False.

The truth table below shows that $p \lor \lnot p$ is a \todo{tautology} and that $p \land \lnot p$ is a \todo{contradiction}.
\begin{center} 
\begin{tabular}{ | c | c | c |c |c |c |c |c |} 
  \hline
 $p$ & T & $ p \lor \lnot p$&F & $p \land \lnot p$\\
 \hline
  T &T  &T&F&F\\ 
  F & T & T&F&F\\ 
  \hline
\end{tabular}
\end{center}
\end{exmp}

\begin{exmp}[Here are some Logical Equivalences]
\begin{align}
p \land T &\equiv p  & \text{Identity Laws}\\
p \lor F &\equiv p  & \\
p \lor T & \equiv T & \text{Domination Laws}\\
p \land F & \equiv F & \\
p \lor p & \equiv p & \text{Idempotent Laws}\\
p \land p & \equiv p & \\
\lnot \lnot p & \equiv p & \text{Double Negation}\\
p \lor q & \equiv q \lor p & \text{Commutative Laws}\\
p \land q & \equiv q \land p & \\
(p \lor q) \lor r & \equiv p \lor (q \lor r) & \text{Associative Laws}\\
(p \land q) \land r & \equiv p \land (q \land r) \\
p \lor (q\land\land r) & \equiv (p \lor q ) \land 
 (p \lor r)  & \text{Distributive Laws}\\
p \land (q\lor r) & \equiv (p \land q) \lor (p \land r) & \\
\lnot (p\land q) & \equiv \lnot p \lor \lnot q & \text{De Morgan's Laws}\\
\lnot (p\lor q) & \equiv \lnot p \land \lnot q & \\
p \lor (p \land q) & \equiv p & \text{Absorption Laws}\\
p \land (p \lor q)& \equiv p &\\
p \to q & \equiv q \lor \lnot p   & \text{$\lor$ restatement of implication}\\
p \to q & \equiv \lnot q \to \lnot p & \text{Contraposition}\\
(p \to q) \land (p \to r) &\equiv p\to (q \land r)& \text{Conjunction of implications}\\
(p \to r) \land (q \to r) &\equiv (p \lor q) \to r & \text{will show this below}  \\ 
(p \to q) \lor (p \to r) &\equiv p\to q \lor r& \text{disjunction of implications}\\
(p \to r) \lor (q \to r) &\equiv (p\land q)\to r\\
p \iff q & \equiv (p\to q)\land(q\to p)  & \text{Conjunction of implications}\\
p \iff q & \equiv \lnot p \iff \lnot q  & \text{Negation restatement}\\
\lnot(p \iff q) & \equiv \lnot p \iff q  & \text{Negation}
\end{align}
\end{exmp}


\newpage
\begin{exmp}[Showing logical equivalence]
Suppose that we want to prove $(p \to r) \land (q \to r) \equiv (p \lor q) \to r$ using only Laws we have previously established.
we may do so by transforming the LHS to the RHS one transformation at a time. Each transformation, we should state which Law we are using for example, this is a long hard one:

\begin{align*}
(p \to r) \land (q \to r) \equiv& (r \lor \lnot p) \land (r\lor \lnot q)\\
&\equiv (r \lor (\lnot p \land \lnot q))\\
&\equiv (r \lor \lnot (p \lor q))\\
&\equiv (p \lor q) \to r
\end{align*}
\end{exmp}

\newpage 
\subsection{Predicates and Quantifiers}

\begin{defn}[predicate]
A \textbf{predicate} is a boolean valued function-- that is, $P$ is a function $P: \mathfrak{U} \to \{True,False\}$
\end{defn}

\begin{defn}[Universal Quantification]
The \textbf{universal quantification} of $P(x)$ is the statement
“$P(x)$ for all values of $x$ in the domain.”
The notation $\forall x P(x)$ denotes the universal quantification of $P(x)$. Here $\forall$ is called the
universal quantifier.We read $\forall xP(x)$ as “for all $x P(x)$” or “for every  $xP(x)$.” An element
for which $P(x)$ is false is called a \textbf{counterexample} of $ \forall xP(x)$.
\end{defn}
\begin{exmp}

Suppose that $\mathfrak{U}:= \{2,3,4,5,6,7,8,9\}$.
Determine the truth of the proposition "All odd numbers in $\mathfrak{U}$ are prime." 
First we define the predicates
$P(x):=\text{$x$ is prime}$, 
$S(x):= \text{$x$ is odd}$.

Using Universal quantification, the expression above is equivalent to $\forall x (S(x)\to P(x))$.

This is equivalent to the statement:
$(S(1) \to P(1)) \land (S(2) \to P(2)) \land (S(3) \to P(3)) \cdots$

which is false if $S(i)\to P(i)$ is false for some i in $\mathfrak{U}$.

\begin{center} 
\begin{tabular}{ | c | c | c |c |c |c |c |c |} 
  \hline
 $x$& $P(x)$ & $S(x)$ & $S(x) \to P(x)$ \\
 \hline
  2 & T & F& T\\ 
  3 & T & T & T\\ 
  4 & F & F & T\\ 
  5 & T & T & T\\ 
  6 & F & F & T\\ 
  7 & T & T & T\\ 
  8 & F & F & T\\ 
  9 & F & T & F\\ 
  \hline
\end{tabular}

\end{center}

\vspace{2cm}
\end{exmp}
\begin{defn}[Existential Quantification]
The \textbf{existential quantification} of $P(x)$ is the proposition
“There exists an element $x$ in the domain such that $P(x)$.”
We use the notation $\exists xP(x)$ for the existential quantification of $P(x)$. Here $\exists$ is called the
existential quantifier.
\end{defn}
\begin{exmp}[Existential]
Let $P(x)$ denote the statement “$x^2 = 8 $.” What is the truth value of the quantification $\exists xP(x)$,
where the domain is all integers?
\begin{proof}
$\exists x P(x)$ is True if there exists an integer such that $x^2=8$. Since no such integer exists, we $\exists x P(x)$ is False.
\end{proof}
\end{exmp}
\begin{defn}[set notation] \ \\
The predicate
\\
$P(x):= \text{"x is in the set $P$"}$ is so common that we have special notation for it, $x \in P$.
\\
We may restate statements in predicate logic using this alternate notation:
\\
$\forall x (P(x) \implies R(x))\iff \forall x \in P, R(x)$
\\
We read this as "For all $x$ in $P$, $R(x)$". With the meaning, "For all $x$ in $P$, $R(x)$ is True."
\\
$\exists x (P(x) \land  R(x))\iff \exists x \in P s.t. R(x)\iff \exists x \in P | R(x)$
\\
We read this as "There exists an $x$ in $P$, such that $R(x)$". With the meaning, "There exists an $x$ in $P$, such that $R(x)$ is True."
\\
We read the statement $x\notin P$, as "$x$ not in $P$".
\end{defn}
\begin{rem}
These statements are not a part of the formal language, but they are great at making arguments  more clear to the reader. (Another example of this is when we write $\cdots$. While this is not in the formal language, using $\cdots$ can help make our arguments easier to follow and we like helping the reader whenever we can.)
\end{rem}
\begin{exmp}[translation]
Translate the following proposition to predicates, quantifiers
"There is a real number which is not rational, but it's square is a rational"

\begin{proof}
We know that the proposition is equivalent to the statement "There is a real number which is not rational, and it's square is a rational".
\\
Recall, the statement"$x$ is real" is equivalent to $x \in \mathbb{R}$.
\\
Recall, the statement"$x$ is rational" is equivalent to $x \in \mathbb{Q}$.
\\
Putting this together, we have
$\exists x \in \mathbb{R} |( x \notin \mathbb{Q} \land x^2 \in \mathbb{Q} )$
\end{proof}
\begin{proof}
We may also write this statement without set inclusion notation.



\end{proof}



\end{exmp}

\begin{exmp}
Translate each of these statements into logical expressions
using predicates, quantifiers, and logical connectives.
\begin{itemize}
\item a) No one is perfect.
\begin{proof}
$\mathfrak{U}:=\{\text{all people}\}$ 
\newline
$P(x):=\text{"x is perfect"}$
\newline
"It is not that case that someone is perfect"
\newline
"It is not that (there exists a person who is perfect)"
\newline
$\lnot(\exists x, P(x)) $
\end{proof}
\item b) Not everyone is perfect.
\begin{proof}
$\mathfrak{U}:=\{\text{all people}\}$ 
\newline
$P(x):=\text{"x is perfect"}$
\newline
" It is not the case that (For all people, people are perfect)"
\newline
$\lnot(\forall x, P(x)) $
\end{proof}
\item c) All your friends are perfect.
\begin{proof}
$\mathfrak{U}:=\{\text{all people}\}$ 
\newline
$F(x):=\text{"x is my friend"}$
\newline
$P(x):=\text{"x is perfect"}$
\newline
"For all people, if someone is your friend, then they are perfect."
\newline
$\forall x ( F(x) \to P(x) ) $
\end{proof}
\item d) At least one of your friends is perfect.
\begin{proof}
$\mathfrak{U}:=\{\text{all people}\}$ 
\newline
$F(x):=\text{"x is my friend"}$
\newline
$P(x):=\text{"x is perfect"}$
\newline
version 1:
\newline
"There exists a person, who is your friend and who is perfect."
\newline
$\exists x (F(x) \land P(x)) $
\newline
Doesn't WORK
\newline
version 2:"There exists a person, who if they are your friend, they are perfect."
\newline
$\exists x (F(x) \to P(x)) $
\end{proof}
\end{itemize}
\end{exmp}
$\mathfrak{U}:=\{\text{Unicorns}\}$
\newline
$F\equiv\exists x \lnot (P(x))\equiv \lnot \forall x P(x)$
\newline
$F\equiv\lnot \forall x P(x)$
\newline
$T\equiv \forall x P(x)$


\begin{defn}[Demorgan's Law for quantification]

\[\lnot (\forall x P(x)) \equiv \exists x \lnot( P(x))\]
\[\lnot (\exists x P(x)) \equiv \forall x \lnot( P(x))\]

\end{defn}
\begin{exmp}

What are the negations of the following statements

"Everyone is a cop or a robber"

"Everyone who loves cats loves dogs"
\end{exmp}
\begin{proof}
Let 
\\$C(x):=\text{$x$ is a cop }$ and
\\$R(x):=\text{$x$ is a robber}$.
\\Then "Everyone is a cop or a robber" may be written as $\forall x (C(x) \lor R(x))$.
The negation of which is $\lnot \forall x (C(x) \lor R(x))$.
\\ This is equivalent to $\exists x \lnot (C(x) \lor R(x))$, which may be written as $\exists x (\lnot C(x) \land \lnot R(x))$.
\\ putting this back in to plain language gives us the statement
\\"There exists someone who is neither a cop nor a robber"
\end{proof}
"Everyone who loves cats loves dogs"
\begin{proof}
Let 
\\$C(x):=\text{$x$ loves cats}$ and
\\$D(x):=\text{$x$ loves dogs}$.
\\Then "Everyone who loves cats loves dogs" may be written as $\forall x (C(x) \to D(x))$.
The negation of which is $\lnot \forall x (C(x) \to R(x))$.
\\ This is equivalent to $\exists x \lnot (C(x) \to D(x))$, which may be written as $\exists x \lnot ( D(x) \lor \lnot C(x) )$, which then may be rewritten as $\exists x ( \lnot D(x) \land C(x) )$.
\\ putting this back in to plain language gives us the statement
\\"There exists someone who dislikes dogs and who likes cats"
\end{proof}

\subsection{Nested Quantifiers}

\begin{exmp}
Are the following equivalences?

\begin{align}
\forall x (\forall y P(x,y))\equiv & \forall y (\forall x P(x,y)) \text{ this is true!}\\
\forall x (\exists y P(x,y))\equiv & \exists y (\forall x P(x,y)) \text{ this is false!}\\
\exists x (\exists y P(x,y))\equiv & \exists y (\exists x P(x,y)) \text{ this is true!}
\end{align}
\end{exmp}
\begin{exmp}
Suppose that $P(x,y)$ uses $\mathfrak{U}=:\{1,2\}$ how would I state the following using universal and existential quantifiers?


\begin{align}
(P(1,1)\land P(1,2))\lor(P(2,1) \land P(2,2))\equiv & \exists x \forall y P(x,y)\\
(P(1,1)\lor P(1,2))\lor(P(2,1) \lor P(2,2))\equiv & \exists x \exists y P(x,y)\\
(P(1,1)\land P(1,2))\land(P(2,1) \land P(2,2))\equiv &\forall x \forall y P(x,y)\\
(P(1,1)\lor P(1,2))\land(P(2,1) \lor P(2,2))\equiv &\forall x \exists y P(x,y)
\end{align}
\end{exmp}
\newpage
\begin{exmp}[]\ \\
Rewrite the following expressions into equivalent expressions where all negation operators immediately precede predicates. Show all intermediate steps. For example, $P(x) \land \lnot Q(x)$ is an acceptable form, $\lnot(P(x) \land Q(x))$ is not acceptable.

    \[\lnot \exists x \exists y \bigl( \lnot S(x) \land \lnot T(y) \bigr)\]
    \vspace{3cm}
%Solution:
\begin{proof}
 \begin{align}
\lnot \exists x \exists y \bigl( \lnot S(x) \land \lnot T(y) \bigr)&\equiv \forall x \lnot \exists y \bigl( \lnot S(x) \land \lnot T(y) \bigr)\text{DeMorgans}\\
&\equiv \forall x \forall y \lnot  \bigl( \lnot S(x) \land \lnot T(y) \bigr)&\text{DeMorgans}\\
&\equiv \forall x \forall y  \bigl( \lnot \lnot S(x) \lor \lnot \lnot T(y) \bigr)&\text{DeMorgans}\\
&\equiv \forall x \forall y  \bigl( S(x) \lor T(y) \bigr)&\text{DeMorgans}
 \end{align}
\end{proof}
\end{exmp}
\begin{exmp}[]\ \\
Rewrite the following expressions into equivalent expressions where all negation operators immediately precede predicates. Show all intermediate steps. For example, $P(x) \land \lnot Q(x)$ is an acceptable form, $\lnot(P(x) \land Q(x))$ is not acceptable.

    \[\lnot \forall x \bigl( \lnot P(x) \land \lnot \exists y T(y) \bigr)\]
    \vspace{3cm}
%Solution:
\begin{proof}
 \begin{align}
\lnot \forall x \bigl( \lnot P(x) \land \lnot \exists y T(y) \bigr)&\equiv  \exists x \lnot \bigl( \lnot P(x) \land \lnot \exists y T(y) \bigr)&\text{DeMorgans negation of universal quantifier}\\
&\equiv  \exists x  \bigl( \lnot(\lnot P(x)) \lor \lnot(\lnot \exists y T(y) )\bigr)&\text{DeMorgans}\\
&\equiv  \exists x  \bigl( P(x)) \lor \exists y T(y) \bigr)&\text{DeMorgans}\\
 \end{align}
\end{proof}
\end{exmp}
\begin{exmp}
Translate the following arguments into logical statements using predicates, logical connectives, and quantifiers as appropriate.

  \begin{enumerate}
      \item
    All dogs are mammals.\\
    Some dogs go to heaven.\\
    Some mammals go to heaven
%%%%%%%%%%%%%%%%% Add your solution below%%%%%%%%%%%%%%%%%%
\vspace{3cm}
\begin{proof}
Define the predicates,
$D(x):=\text{x is a Dog}$\\
$M(x):=\text{x is a mammal}$\\
$H(x):=\text{x goes to heaven}$\\
Then
\begin{align}
\text{All dogs are mammals.}\equiv& \forall x (D(x)\to M(x))\\
\text{Some dogs go to heaven.}\equiv& \exists x (D(x)\land H(x))\\
\text{Some mammals go to heaven}\equiv& \exists x(M(x) \land H(x))
\end{align}


\end{proof}

  \item
    All Cats are great.\\
    Some cats are awful.\\
    Some animals are great and awful.
    \vspace{3cm}
    \begin{proof}
Define the predicates,\hspace*{\fill} \\ 
$C(x):=\text{x is a Cat}$ \hspace*{\fill} \\ 
$G(x):=\text{x is Great}$\hspace*{\fill} \\ 
$A(x):=\text{x is an animal}$\hspace*{\fill} \\ 

Then
\begin{align}
\text{All Cats are great.}\equiv& \forall x (C(x)\to G(x))\\
\text{Some cats are awful.}\equiv& \exists x (C(x)\land A(x))\\
\text{Some animals are great and awful.}\equiv& \exists x (A(x) \land G(x))
\end{align}


\end{proof}
    
  \end{enumerate}
%%%%%%%%%%%%%%%%% Add your solution below%%%%%%%%%%%%%%%%%%
\begin{align}
\text{this is stuff}
\end{align}
\end{exmp}

\begin{exmp}
Express each of the following using mathematical and logical operators,
predicates, and quantifiers. Define the domain of each variable used.
(a) Any rational number can be expressed as the quotient of two integers.
(b) There are no real numbers x and y such that x = y/0.

\end{exmp}
\vspace{3cm}

\newpage
\subsection{Rules of Inference}
\begin{defn}[]
An \textbf{argument} in propositional logic is a sequence of propositions. All but the final proposition
in the argument are called \textbf{premises} and the final proposition is called the \textbf{conclusion}. An
argument is \textbf{valid} if the truth of all its premises implies that the conclusion is true.
An argument form in propositional logic is a sequence of compound propositions involving
propositional variables. An argument form is valid no matter which particular propositions
are substituted for the propositional variables in its premises, the conclusion is true if
the premises are all true.
\end{defn}
\begin{defn}[rules of inference]
\textbf{Rules of inference} can be used as building blocks to construct more complicated valid argument forms.
We will now introduce the most important rules of inference in propositional logic.
\end{defn}
This notation is incredibly cumbersome. Instead of writing things in this way:
\begin{tabular}{cl}
    & $A$\\
    & $B$\\
    \hline
    $\therefore$ & $C$
  \end{tabular}
I will be using this
$A,B \vdash C$.

\begin{defn}[Modus ponens]
$p,p\to q \vdash q$
\end{defn}

\begin{defn}[Modus tollens]
$\lnot q,p\to q \vdash \lnot p$
\end{defn}
\begin{defn}[Hypothetical syllogism]
$(p \to q) , (q \to r) \vdash (p \to r)$
\end{defn}
\begin{defn}[Disjunctive syllogism]
$(p \lor q),\lnot p \vdash q$
\end{defn}
\begin{defn}[Addition]
$p \vdash p \lor q$
\end{defn}
\begin{defn}[Simplification]
$p \land q\vdash p$
\end{defn}
\begin{defn}[Conjunction]
$p,q \vdash p \land q$
\end{defn}
\begin{defn}[Resolution]
$p \lor q, \lnot p \lor r  \vdash q \lor r$
\end{defn}
\newpage 
\begin{rem}
We can always think of our truth tables (even if we aren't using them explicitly!)\\
Suppose we want to show the following:\\
If $a\vdash b$ and $b \vdash c$, then $a\vdash c$. \\
This should remind you of the hypothetical syllogism\\
We can say this as "If you can prove the conclusion $b$ from the premise $a$, and you can prove the conclusion $c$ from the premise $b$, then you can prove the conclusion $c$ from the premise $a$."
\\
This means that propositions we have established using laws of inference may be used.
\\
This should make sense with our model interpretation:\\
"If in every world in which a is true, b is also true" and "If every world in which b is true, c is true" then "every world in which a is true, c is true". 
\\
\end{rem}
\begin{exmp}
Suppose that \\
$a :=p\land (\lnot q \land \lnot r)$,\\
$b:= (\lnot q \land \lnot r)$, and \\
$c:=\lnot q \to \lnot r$.\\
Show that
$a\vdash c$.
\begin{proof}
Lets restrict our attention to worlds where $a$ is true. Below are all the assignments in which $a$ is true.\\
\begin{tabular}{ | c | c | c |c |c |} 
  \hline
 $p$ $q$ $r$ & premise & $p \land (\lnot q \land \lnot r)$& conclusion & $\lnot q \land \lnot r$\\
 \hline
  T F F  &&T &&T\\ 
 \hline
\end{tabular}\\
We notice that in all of these worlds (assignments) $b$ is also true.\\
We now restrict our attention only to worlds (assignments) where $b$ is true.\\
Below are all the assignments in which $b$ is true.\\
\begin{tabular}{ | c | c | c |c |c |} 
  \hline
 $p$ $q$ $r$ & premise & $\lnot q \land \lnot r$& conclusion & $\lnot q \to \lnot r$\\
 \hline
  T F F  &&T &&T\\
  F F F  &&T &&T\\ 
 \hline
\end{tabular}\\
We notice that in all of these worlds (assignments) $c$ is also true.\\
Since every assignment in which $a$ is true, $b$ is true, and in every assignment in which $b$ is true, $c$ is true, every assignment in which $a$ is true, $c$ is true.\\
We can think of constructing an argument in this way, by restricting our attention to only some models (where it can be easier to see )\\

\begin{tabular}{ | c | c | c |c |c |} 
  \hline
 $p$ $q$ $r$ & $p \land (\lnot q \land \lnot r)$&$(\lnot q \land \lnot r)$&$\lnot q \to \lnot r$\\
 \hline
  T T T  &F &F&\textbf{T}\\
  T T F  &F &F&\textbf{T}\\
  T F T  &F &F&F\\
  F T T  &F &F&\textbf{T}\\
  T F F  &\textbf{T} &\textbf{T}&\textbf{T}\\ 
  F T F  &F &F&\textbf{T}\\ 
  F F T  &F &F&F\\ 
  F F F  &F &\textbf{T}&\textbf{T}\\ 
 \hline
\end{tabular}\\
Any assignment where the first column is true, the second column is true, and any assignment that the second column is true the third column is true, so every assignment where the first column is true, the last column is true.
\end{proof}
\end{exmp}

\newpage

\begin{exmp}[] \ \\

\begin{enumerate}
\item (this is a formal argument form, you should be able to do this)\\
Show\\
\begin{tabular}{cl}
    & $p \to q$\\
    & $\lnot p \to r$\\
    & $r \to s $\\
    \hline
    $\therefore$ & $\lnot q \to s$
  \end{tabular}\\

\renewcommand{\arraystretch}{2}

\begin{tabularx}{\textwidth - \leftmargin}{clX}
& Argument Steps: & Reason/Rule/Law of Logic:\\

a) & $p\to q$ & \text{ Premise}
\\
\cline{3-3}
b) & $\lnot q \to \lnot p$  & \text{ contraposition of a)}
\\
\cline{3-3}

c) & $\lnot p \to r$ & \text{ Premise}
\\
\cline{3-3}

d) &  $\lnot q \to r$ & \text{ Hypothetical Syllogism of b) and c)}
\\
\cline{3-3}
e) & $r \to s$ & \text{ Premise}
\\
\cline{3-3}
f) & $\lnot q \to s$ & \text{ Hypothetical Syllogism of d) and e)}
\\
\cline{3-3}
\end{tabularx}
\renewcommand{\arraystretch}{1}

\item (This is the equivalent informal argument form, this is how most people will prove things)\\
These should all be written neatly in full sentences. 

Show $p \to q, \lnot p \to r,r \to s \vdash \lnot q \to s$.
\begin{proof}
By contraposition, $p\to q\vdash \lnot q \to \lnot p $.\\
Because we have proven $\lnot q \to \lnot p$ using our premises, we may now use it (as though it was a premise).\\
By hypothetical syllogism, $\lnot q \to \lnot p, \lnot p \to r \vdash  \lnot q \to r$.\\
By hypothetical syllogism, $\lnot q \to r, r \to s \vdash  \lnot q \to s$.\\
\end{proof}
\end{enumerate}
\end{exmp}
\vspace{4cm}
\newpage 
\begin{defn}[Universal Instantiation]
$\forall x P(x) \vdash P(c)$
\end{defn}
\begin{defn}[Universal Generalization]
$\text{for arbitrary c } P(c)  \vdash \forall x P(x) $
\end{defn}
\begin{defn}[Existential Instantiation]
$\exists x P(x) \vdash \text{for some } c, P(c)$
\end{defn}
\begin{defn}[Existential Generalization]
$ P(c) \text{ for some c } \vdash \exists x P(x)$
\end{defn}
\begin{exmp}
Show,
$\exists x(C(x)\land \lnot B(x)), \forall x(C(x) \to P(x)) \vdash \exists x(P(x)\land \lnot B(x))$
\\
(I will do this informally for the sake of brevity, but we should be able to construct the formal argument form also).
\begin{proof}

\end{proof}
By existential instantiation,$\exists x(C(x)\land \lnot B(x))\vdash$ for some $d$ $(C(d)\land \lnot B(d))$. \\
By simplification, from $(C(d)\land \lnot B(d))$ we have $ C(d)$.\\
By universal modus ponens,from 
$\forall x(C(x) \to P(x))$ and $C(d)$ we know $P(d)$.\\
Using, simplification of $C(d)\land \lnot B(d)$ we have  $\lnot B(d)$. 
\\
Conjunction of $\lnot B(d)$, $P(d)$ yields, $P(d) \land \lnot B(d)$.
\\ Thus, through Existential Generalization, we see $\exists x(P(x)\land \lnot B(x))$





\end{exmp}


\begin{exmp}
Determine the truth value of the statement $\exists x \forall y (x\leq y^2)$
if the domain for the variables consists of
\newline
\begin{enumerate}
\item a) the positive real numbers.
\vspace{3cm}
\begin{proof}
We will prove that this statement is false by proving the truth of it's negation.
We begin by rewriting the negation:
\begin{align}
\lnot(\exists x \forall y (x\leq y^2))&\equiv \forall x \lnot \forall y (x\leq y^2))\\
&\equiv \forall x  \exists y \lnot  (x\leq y^2))\\
&\equiv \forall x  \exists y (x > y^2)).
\end{align}
To show the truth of this equivalent statement we will show that for all x in the domain, there exists a y in the domain such that $x>y^2$.
\\
Choose an arbitrary $c$ in the positive real numbers. Define $d:=\sqrt{c/2}$. Notice that $d$ is a positive real number.
\\
It is clear that for arbitrary $c$, for some $d$, $c>d^2$.\\
 for arbitrary $c\in \mathbb{R}$, for some $d\in \mathbb{R}$,$(c > d^2) $.\\
By the inference rule of Existential Generalization,
for arbitrary $c\in \mathbb{R}$, $\exists y \in \mathbb{R}$ such that $(c > y^2) $.\\
By the inference rule of Universal Generalization,
$\forall x, \exists y (x > y^2)$.\\
Therefore, $\lnot(\exists x \forall y (x\leq y^2))$.
\end{proof}
\newpage 
\item b) the integers.
\vspace{3cm}
\begin{proof}
We begin by rewriting the proposition,
\begin{align}
\exists x \forall y (x\leq y^2))\equiv& \text{for some } c,\forall y (c\leq y^2))\text{ Universal Instantiation}\\
\equiv& \text{for some } c,\text{for all d,} (c\leq d^2))\text{ Universal Instantiation}.
\end{align}
Let $c=0$ and $d\in \mathbb{Z}$. Then $\abs{d}\geq 0$ and $d^2\geq 0$.\\
Therefore, for $c=0$,$\forall d\in \mathbb{Z}$, $ (c\leq d^2))$.\\
By Universal Generalization, $c=0$,$\forall y$, $ (c\leq y^2))$.\\
By Existential Generalization, $\exists x$$\forall y$, $ (x\leq y^2))$.\\

\end{proof}
\noindent
\item c) the nonzero real numbers.
\vspace{3cm}
\newline
\end{enumerate}
\end{exmp}
\vspace{3cm}

\begin{exmp}{Importance of Formalizing arguments}
(Because Tony wanted a paradox and I feel bad for making him type in .tex, here is Curry's Paradox)\\

Analyze the truth of the following sentence.\\

"If this proposition is true, then pigs can fly"\\
\end{exmp}
\begin{proof}
$P:=$"If this proposition is true, then then pigs can fly"\\
$Q:=$"pigs can fly"\\
$P:=P\to Q$
If $P$ then by modus ponens, $Q$. this is absurd so $P$ is false. Since P is false, $P\to Q$ is True, $P$ shares a truth value with $P\to Q$, so $P$ is True, therefore $Q$.
\end{proof}



\newpage

\section{Sets, Functions, Sequences, Sums, and Matrices}
\subsection{Sets}

\begin{defn}[Set]
A set is an unordered collection of objects, called elements or members of the set. A set is
said to contain its elements.We write $a \in A$ to denote that $a$ is an element of the set $A$. The
notation $a \notin A$ denotes that $a$ is not an element of the set $A$.
\end{defn}
\begin{defn}[Equality]
Two sets are equal if and only if they have the same elements. Therefore, if $A$ and $B$ are sets,
then $A$ and $B$ are equal if and only if $\forall x(x \in A \leftrightarrow x \in B)$.We write $A = B$ if A and B are
equal sets.
\end{defn}
\begin{defn}[Subset]
The set $A$ is a subset of $B$ if and only if every element of $A$ is also an element of $B$. We use
the notation $A \subseteq B$ to indicate that $A$ is a subset of the set $B$.
\end{defn}
\begin{defn}[empty set]
The set with no element is called the empty set and denoted $\emptyset$ or $\{\}$. We use the language 'the' empty set and not 'an' empty set because all empty sets are equal to one another.
\end{defn}
\begin{thm}[trivial subsets]
For every set $S$, $\emptyset \subseteq S$ and $S \subseteq S$.
\end{thm}

\begin{defn}[cardinality]
Let $S$ be a set. If there are exactly $n$ distinct elements in $S$ where $n$ is a nonnegative integer,
we say that $S$ is a finite set and that $n$ is the cardinality of S. The cardinality of $S$ is denoted
by $\card{S}$.
\end{defn}
\begin{defn}[infinite]
A set is said to be infinite if it is not finite.
\end{defn}
\begin{defn}[Power Set]
Given a set $S$, the power set of $S$ is the set of all subsets of the set $S$. The power set of $S$ is
denoted by $\mathcal{P}(S)$.
\end{defn}
\begin{defn}[$n$-tuple]
The ordered $n$-tuple $(a_1, a_2, . . . , a_n)$ is the ordered collection that has $a_1$ as its first element,
$a_2$ as its second element,$\cdots$, and an as its $n$th element.
\end{defn}
\begin{defn}[Cartesian Product]
Let $A$ and $B$ be sets. The Cartesian product of $A$ and $B$, denoted by $A \times B$, is the set of all
ordered pairs $(a, b)$, where $a \in A$ and $b \in B$. Hence,
$A \times B = {(a, b) | a \in A \land b \in B}$.
\end{defn}
\begin{defn}[Cartesian Product of Sets]
The Cartesian product of the sets $A_1,A_2, \cdots , A_n,$ denoted by $A_1 \times A2 \times \cdots \times A_n$, is the
set of ordered $n$-tuples $(a_1, a_2, \cdots , a_n)$, where $a_i$ belongs to $A_i$ for $i = 1, 2, \cdot , n$. In other
words,
\[A_1 \times A_2 \times \cdots \times A_n = \{(a_1, a_2, \cdots , a_n) | a_i \in A_i \text{ for }i = 1, 2,\cdots, n\}\].
\end{defn}
\begin{exmp}
Draw out $\mathcal{P}(\{a,b,c\})$ and draw an arrow from the set  $A$ to the set $B$ if $A\subseteq B$ and $\card{B}-\card{A}=1$.\\
\vspace{2cm}\\
Describe $A\cap B$ using directed paths to $A$ and $B$.\\
\vspace{1cm}\\
Describe $A\cup B$ using directed paths from $A$ and $B$.\\
\vspace{1cm}\\

\end{exmp}

\newpage
\subsection{Set Operations}
\begin{defn}[union]
Let $A$ and $B$ be sets. The union of the sets $A$ and $B$, denoted by $A \cup B$, is the set that contains
those elements that are either in $A$ or in $B$, or in both.
\end{defn}
\begin{defn}[intersection]
Let $A$ and $B$ be sets. The intersection of the sets $A$ and $B$, denoted by $A \cap B$, is the set
containing those elements in both $A$ and $B$.
\end{defn}
\begin{defn}[disjoint]
Two sets are called disjoint if their intersection is the empty set.
\end{defn}
\begin{defn}[difference]
Let $A$ and $B$ be sets. The difference of $A$ and $B$, denoted by $A - B$ (or more commonly $A \setminus B$), is the set containing those
elements that are in $A$ but not in $B$. The difference of $A$ and $B$ is also called the complement
of $B$ with respect to $A$.
\end{defn}
\begin{defn}[compliment]
Let $U$ be the universal set. The complement of the set $A$, denoted by $\overline{A}$, is the complement
of $A$ with respect to $U$. Therefore, the complement of the set $A$ is $U - A$ ($U \setminus A$).
\end{defn}
\begin{exmp}Deriving Set Identities
\begin{align}
p \land T &\equiv p  &A\cap U=A \hspace{1cm}& \text{Identity Laws}\\
p \lor F &\equiv p  &A\cup \emptyset=A \hspace{1cm}& \\
p \lor T & \equiv T &A\cup U=U \hspace{1cm}& \text{Domination Laws}\\
p \land F & \equiv F &A\cap \emptyset=\emptyset \hspace{1cm}& \\
p \lor p & \equiv p &A\cup A=A \hspace{1cm}& \text{Idempotent Laws}\\
p \land p & \equiv p &A\cap A=A\hspace{1cm}& \\
\lnot \lnot p & \equiv p &A=((A)^c)^c\hspace{1cm}& \text{Complementation}\\
p \lor q & \equiv q \lor p &A\cap B=B\cap A\hspace{1cm}& \text{Commutative Laws}\\
p \land p & \equiv q \land p &A\cup B=B\cup A\hspace{1cm}& \\
(p \lor q) \lor r & \equiv p \lor (q \lor r) &(A\cup B)\cup C=A\cup (B\cup C)\hspace{1cm}& \text{Associative Laws}\\
(p \land q) \land r & \equiv p \land (q \land r)&(A\cap B)\cap C=(A\cap B)\cap C\hspace{1cm}& \\
p \lor (q\land r) & \equiv (p \lor q ) \land 
 (p \lor r)  &A \cup (B\cap C) = (A \cup B ) \cap 
 (A \cup C)\hspace{1cm}& \text{Distributive Laws}\\
p \land (q\lor r) & \equiv (p \land q) \lor (p \land r) &\underline{\hspace{5cm}}\hspace{1cm}& \\
\lnot (p\land q) & \equiv \lnot p \lor \lnot q & \overline{A\cap B}= \overline{A} \cup \overline{B} \hspace{1cm}& \text{De Morgan's Laws}\\
\lnot (p\lor q) & \equiv \lnot p \land \lnot q &\underline{\hspace{5cm}}\hspace{1cm}& \\
p \lor (p \land q) & \equiv p &\underline{\hspace{5cm}}\hspace{1cm}& \text{Absorption Laws}\\
p \land (p \lor q)& \equiv p &\underline{\hspace{5cm}}\hspace{1cm}&\\
p \lor \lnot p= T  &\hspace{3cm}&\underline{\hspace{5cm}}\hspace{1cm}& \text{Complement Laws}\\
p \land \lnot p=F & \equiv q \lor \lnot p   &\underline{\hspace{5cm}}\hspace{1cm}&
\end{align}
\end{exmp}
\begin{defn}[arbitrary union]
The union of a collection of sets is the set that contains those elements that are members of
at least one set in the collection.\\
We use the notation
$A_1\cup A_2 \cup \cdots \cup A_n = \union_{i=1}^{n}A_n$
to denote the union of the sets $A_1,A_2, \cdots, A_n$.
\end{defn}
\begin{defn}[arbitrary intersection]
The intersection of a collection of sets is the set that contains those elements that are members
of all the sets in the collection.\\
We use the notation
$A_1\cap A_2 \cap \cdots \cap A_n = \intersect_{i=1}^{n}A_n$
to denote the intersection of the sets $A_1,A_2, \cdots, A_n$.
\end{defn}


\begin{exmp}
Suppose that $A$ and $B$ are subsets of the universal set $U$. 
\\
Match the sets below with equivalent sets from the bank below.
\begin{align}
&\{x| x\in A \land x \notin B\} \hspace{3cm}&\underline{\hspace{1cm}
%YourSolutionHere
}\\
&\{x| (x\in A) \to (x \in B)\} \hspace{3cm}&\underline{\hspace{1cm}
%YourSolutionHere
}\\
&\{x\in A| (x\in A) \to (x \in B)\} \hspace{3cm}&\underline{\hspace{1cm}
%YourSolutionHere
}\\
&\{x| (x\in A) \leftrightarrow (x \in B)\} \hspace{3cm}&\underline{\hspace{1cm}
%YourSolutionHere
}\\
&\{x| (x\notin A) \lor (x \notin B)\} \hspace{3cm}&\underline{\hspace{1cm}
%YourSolutionHere
}
\end{align}
\end{exmp}
\begin{exmp}
Let $A$,$B$ and $C$ be subsets of $U$.\\
If $A\cap C=B\cap C$ and $A\cup C =B\cup C$ then $A=B$.
\end{exmp}
\begin{proof}
We begin by showing that $A\subseteq B$.\\
Suppose $x\in A$.\\
There are two cases\\
Case 1: $x\in C$\\
Suppose $x \in C$. Then $x\in A\cap C$ and since $A\cap C=B\cap C$, $x\in B\cap C$, therefore $x \in B$.\\
Case 2: $x\notin C$\\
Suppose $x\notin C$, 
then $x$ in $A\cup C$. 
Since $A\cup C=B\cup C$, 
$x\in B\cup C$. 
Therefore $x\in B$ or $x \in C$.
Since $x \notin C$ by assumption, $x\in B$.
Since in all cases $x \in B$, we have $x \in A \to x \in B$ and so $A \subseteq B$.
\end{proof}

\begin{lem}
If $A=B\cap A$ and $A =B\cup A$ then $A=B$
\end{lem}
\begin{proof}
\vspace{5cm}
\end{proof}
\begin{exmp}
Show $(A\cup B)\cap (B\cup C)\cap (C\cup A) =(A\cap B)\cup (B\cap C)\cup (C\cap A)$
\end{exmp}
\begin{proof}

\end{proof}

\newpage
\subsection{Functions}

\begin{defn}[function]
Let $A$ and $B$ be nonempty sets. A function $f$ from $A$ to $B$ is an assignment of exactly one
element of $B$ to each element of $A$. We write $f (a) = b$ if $b$ is the unique element of $B$
assigned by the function $f$ to the element $a$ of $A$. If $f$ is a function from $A$ to $B$, we write
$f : A \to B$.
\end{defn}
\begin{defn}[domain codomain]
If $f$ is a function from $A$ to $B$, we say that $A$ is the domain of $f$ and $B$ is the codomain of $f$.
If $f (a) = b$, we say that $b$ is the image of $a$ and $a$ is a preimage of $b$. The range, or image,
of $f$ is the set of all images of elements of $A$. Also, if $f$ is a function from $A$ to $B$, we say
that $f$ maps $A$ to $B$.
\end{defn}
\begin{defn}[function addition and multiplication]
Let $f_1$ and $f_2$ be functions from $A$ to $\mathbb{R}$. Then $f_1 + f_2$ and $f_1f_2$ are also functions from $A$
to $\mathbb{R}$ defined for all $x \in A$ by
$(f_1 + f_2)(x) = f_1(x) + f_2(x)$,
$(f_1\cdot f_2)(x) = f_1(x)f_2(x)$.
\end{defn}
\begin{defn}[image]
Let $f$ be a function from $A$ to $B$ and let $S$ be a subset of $A$. The image of $S$ under the function
$f$ is the subset of $B$ that consists of the images of the elements of $S$.We denote the image of
$S$ by $f (S)$, so
$f (S) = \{t | \exists \in S (t = f (s))\}$.
We also use the shorthand $\{f (s) | s \in S\}$ to denote this set.
\end{defn}
\begin{defn}[injective]
A function $f$ is said to be one-to-one, or an injunction, if and only if $f (a) = f (b)$ implies that
$a = b$ for all $a$ and $b$ in the domain of $f$. A function is said to be injective if it is one-to-one.
\end{defn}

\begin{defn}[onto]
A function $f$ from $A$ to $B$ is called onto, or a surjection, if and only if for every element
$b \in B$ there is an element $a \in A$ with $f (a) = b$. A function $f$ is called surjective if it is onto.
\end{defn}
\begin{defn}[bijection]
The function $f$ is a one-to-one correspondence, or a bijection, if it is both one-to-one and
onto. We also say that such a function is bijective.
\end{defn}

\begin{defn}[inverse]
Let $f$ be a one-to-one correspondence from the set $A$ to the set $B$. The inverse function of
$f$ is the function that assigns to an element $b$ belonging to $B$ the unique element $a$ in $A$
such that $f(a)=b$.

The inverse function of $f$ is denoted by $f^{-1}$. 

Hence, $f^{-1}(b)=a$
when
$f(a)=b$.


\end{defn}

\begin{defn}[composition]
Let $g$ be a function from the set $A$ to the set $B$ and let $f$ be a function from the set $B$ to the
set $C$. The composition of the functions $f$ and $g$, denoted for all $a \in A$ by $f \circ g$, is defined
by
$(f \circ g)(a) = f (g(a))$.
\end{defn}
\begin{defn}[floor]
The floor function assigns to the real number $x$ the largest integer that is less than or equal to
$x$. The value of the floor function at $x$ is denoted by $\lfloor x \rfloor$. The ceiling function assigns to the
real number $x$ the smallest integer that is greater than or equal to $x$. The value of the ceiling
function at $x$ is denoted by $\lceil x \rceil$.
\end{defn}

\newpage
\begin{exmp}{Domain and Range}
Find the Domain and Range of these Functions. 

\begin{itemize}
\item a function which takes in a word and outputs a word that rhymes
\newline
$A:=\{words \}$

$B:=\{ words\}$

$f:A\to B$

\item a function which takes in a bit string and returns the number of ones.

$A:=\{"bit strings" \}$

$B:=\{0, 1, 2, 3, \cdots \}$

$f:A\to B$

\item a function which takes in a group of people and returns the oldest person

$\mathcal{P}(A):=\{S: S\subseteq A\}$

$A:=\mathcal{P}(\{"all people"\})$

$B:=\{\text{"thing"} | \text{"thing is a person"}\}$

$f:A\to B$

\item a function which takes in a pizza topping and returns the group of people who like the pizza topping

$A:=\{\text{"all pizza toppings"}\}$

$B:=\mathcal{P}("all people")$

$f:A\to B$

$f('pepperoni')=\{hussain, will, alex \}$

$\mathcal{P}(A):=\{S |S\subseteq A\}$

$S \in \mathcal{P}(A)$

$S \subseteq A$

\item a function which takes in a set of pizza toppings and returns the group of all people who like these pizza toppings

$A:=\mathcal{P}(\{pizza toppings\})$

$B:=\mathcal{P}(\{all people\})$


\end{itemize}
\end{exmp}
\begin{defn}
We can present a function $f:A\to B$, as a set of tuples $F$, where $F:=\{(a,b)| a\in A, b \in B, f(a)=b\}$
\end{defn}

\begin{exmp}
Determine whether each of these functions from $\{a,b,c,d\}$ to itself is one to one, onto and/or bijective.
\begin{itemize}
\item $F=\{(a,b), (b,a), (c,a),(d,d)\}$

not onto (nothing maps to c)

not one to one (both b and c map to a)
\item $F=\{(a,a), (b,b), (c,c),(d,d)\}$
Since the image of $A$ under $f$ is the codomain, $f$ is onto.

One to one, because $a$ maps to itself.



\end{itemize}
\end{exmp}

\begin{exmp}
Determine whether each of these functions from $\mathbb{R} \to \mathbb{R}$ is one to one, onto and/or bijective.
\begin{itemize}
\item $f(x)=|x-2|+3$

$f(6)=7$ and $f(-2)=7$ but $6\neq -2 $ therefore $f$ is not injective. 


$f(x)>=0+3=3$ so for instance $f(x)\neq -2$ and since $-2$ is in the codomain, $f$ is not onto.

\item $f(x)=x^3$

Suppose $f(a)=f(b)$ then $a^3=b^3$, therefore $a=b$. Therefore injective

$b \in \mathbb{R}$, find an $a$ such that $f(a)=b$ let $a= \sqrt[3]{b}$ so this is onto!

\item $f(x)=2^x$


$\ln(a^b)=b\cdot \ln(a)$


one to one, let $f(a)=f(b)$, then $2^a=2^b$
so 

$\ln(2^a)=\ln(2^b)$

$a \cdot \ln(2)=b \cdot \ln(2)$

$a=b$

It is not onto! we will never find we have negative numbers.

$a^b=a\cdot a\cdot a$ b many times

$a^{-b}=\frac{1}{a^b}$


\end{itemize}
\end{exmp}
\begin{exmp}
Give examples of a function from $\mathbb{N}$ to $\mathbb{N}$ which is.
\begin{itemize}
\item one to one, but not onto

$\mathbb{N}:=\{1,2, \cdots \}$

$f(n)=2\cdot n$

\item onto but not one to one
$\mathbb{N}:=\{1,2, \cdots \}$

$f(n)=\lfloor \frac{n}{2} \rfloor+1$

$f(2)=2$, $f(3)=2$

$f(x)=|4-x|+1$


$3*(n-3)^3+15(n-3)^2$

\item both one to one and onto
$f(n)=n$



\item neither one to one nor onto

$f(n)=1$

$f(n)=n^0$
\vspace{1cm}
\end{itemize}
\end{exmp}
\newpage
\begin{exmp}
Determine whether each of these sets in finite, countably infinite or uncountable
\begin{itemize}
\item the numbers divisible by both 3 and 5
\vspace{1cm}
\item the set of all finite length words of an alphabet with 6 letters
\vspace{1cm}
\item the set of all infinite length words of an alphabet with 6 letters
\vspace{1cm}
\item The interval $(p,q)$ where $p<q$.
\end{itemize}
\end{exmp}
\newpage
\begin{exmp}{Compositions}
\begin{itemize}
\item Show that if $f$ and $g$ are onto then $f \circ g$ is onto. 
\vspace{3cm}
\item Show that if $f$ and $g$ are one to one then $f \circ g$ is one to one. 
\vspace{3cm}
\item Argue that $f$ and $g$ are bijections, then $f \circ g$ is a bijection
\vspace{1cm}
\item Show that if $|A|=|B|$ and $|B|=|C|$ then $|A|=|C|$
\vspace{3cm}
\end{itemize}
\end{exmp}

\newpage

\subsection{Sequence and Summations}
\begin{defn}[sequence]
A sequence is a function from a subset of the set of integers (usually either the set $\{0, 1, 2, \ldots\}$ or the set $\{1, 2, 3, \ldots\}$) to a set $S$. We use the notation $a_n$ to denote the image of the integer $n$. We call $a_n$ a term of the sequence.
\end{defn}

\begin{defn}[geometric progression]
A geometric progression is a sequence of the form $a, ar, ar^2, \ldots, ar^n, \ldots$, where the initial term $a$ and the common ratio $r$ are real numbers.
\end{defn}
\begin{defn}[arithmetic progression]
An arithmetic progression is a sequence of the form $a, a + d, a + 2d, \ldots, a + nd, \ldots$, where the initial term $a$ and the common difference $d$ are real numbers.

\end{defn}
\begin{defn}[Recurrence Relation]
A recurrence relation for the sequence $\{a_n\}$ is an equation that expresses $a_n$ in terms of one or more of the previous terms of the sequence, namely, $a_0, a_1, \ldots, a_{n-1}$, for all integers $n$ with $n \geq n_0$, where $n_0$ is a nonnegative integer. A sequence is called a solution of a recurrence relation if its terms satisfy the recurrence relation.
\end{defn}
\begin{defn}[Fibonacci sequence]
The Fibonacci sequence, $f_0, f_1, f_2, \ldots$, is defined by the initial conditions $f_0 = 0$, $f_1 = 1$, and the recurrence relation
\[f_n = f_{n-1} + f_{n-2}\]
for \(n = 2, 3, 4, \ldots\).
\end{defn}
\begin{defn}[closed form geometric]
If \(a\) and \(r\) are real numbers and \(r \neq 0\), then
\[
\sum_{j=0}^{n} ar^j =
\begin{cases}
    \frac{ar^{n+1} - a}{r - 1} & \text{if } r \neq 1 \\
    (n + 1)a & \text{if } r = 1.
\end{cases}
\]
\end{defn}


\begin{table}[h]
\centering
\caption{Some Useful Summation Formulae.}
\begin{tabular}{|c|c|}
\hline
Sum & Closed Form \\
\hline
$\displaystyle\sum_{k=0}^{n} ar^k$ & $\frac{ar^{n+1} - a}{r - 1}$, $r \neq 1$ \\
\hline
$\displaystyle\sum_{k=1}^{n} k$ & $\frac{n(n + 1)}{2}$ \\
\hline
$\displaystyle\sum_{k=1}^{n} k^2$ & $\frac{n(n + 1)(2n + 1)}{6}$ \\
\hline
$\displaystyle\sum_{k=1}^{n} k^3$ & $\frac{n^2(n + 1)^2}{4}$ \\
\hline
$\displaystyle\sum_{k=0}^{\infty} x^k$, $|x| < 1$ & $\frac{1}{1 - x}$ \\
\hline
$\displaystyle\sum_{k=1}^{\infty} kx^{k-1}$, $|x| < 1$ & $\frac{1}{(1 - x)^2}$ \\
\hline
\end{tabular}
\end{table}
\newpage
\subsubsection{EXAMPLES}
\newpage
\subsection{Cardinality of Sets}
\begin{defn}[Cardinality]
The sets $A$ and $B$ have the same cardinality if and only if there is a one-to-one correspondence from $A$ to $B$. When $A$ and $B$ have the same cardinality, we write $|A| = |B|$.

\end{defn}
\begin{defn}[$|A| \leq |B|$]
If there is a one-to-one function from $A$ to $B$, the cardinality of $A$ is less than or the same as the cardinality of $B$, and we write $|A| \leq |B|$. Moreover, when $|A| \leq |B|$ and $A$ and $B$ have different cardinalities, we say that the cardinality of $A$ is less than the cardinality of $B$, and we write $|A| < |B|$.
\end{defn}

\begin{thm}[Cantor's Diagonalization]

To show that the set of real numbers is uncountable, we suppose that the set of real numbers is countable and arrive at a contradiction. Then, the subset of all real numbers that fall between 0 and 1 would also be countable (because any subset of a countable set is also countable; see Exercise 16). Under this assumption, the real numbers between 0 and 1 can be listed in some order, say, $r_1, r_2, r_3, \ldots$. Let the decimal representation of these real numbers be
\[
\begin{aligned}
    r_1 &= 0.d_{11}d_{12}d_{13}d_{14}\ldots, \\
    r_2 &= 0.d_{21}d_{22}d_{23}d_{24}\ldots, \\
    r_3 &= 0.d_{31}d_{32}d_{33}d_{34}\ldots, \\
    r_4 &= 0.d_{41}d_{42}d_{43}d_{44}\ldots, \\
    \ldots
\end{aligned}
\]
where $d_{ij} \in \{0, 1, 2, 3, 4, 5, 6, 7, 8, 9\}$. (For example, if $r_1 = 0.23794102\ldots$, we have $d_{11} = 2$, $d_{12} = 3$, $d_{13} = 7$, and so on.) Then, form a new real number with decimal expansion
\[
r = 0.d_1d_2d_3d_4\ldots,
\]
where the decimal digits are determined by the following rule:
\[
d_i =
\begin{cases}
    4 & \text{if } d_{ii} = 4, \\
    5 & \text{if } d_{ii} \neq 4.
\end{cases}
\]
(As an example, suppose that $r_1 = 0.23794102\ldots$, $r_2 = 0.44590138\ldots$, $r_3 = 0.09118764\ldots$, $r_4 = 0.80553900\ldots$, and so on. Then, we have $r = 0.d_1d_2d_3d_4\ldots = 0.4544\ldots$, where $d_1 = 4$ because $d_{11} \neq 4$, $d_2 = 5$ because $d_{22} = 4$, $d_3 = 4$ because $d_{33} \neq 4$, $d_4 = 4$ because $d_{44} = 4$, and so on.)

A number with a decimal expansion that terminates has a second decimal expansion ending with an infinite sequence of 9s because $1 = 0.999\ldots$. Every real number has a unique decimal expansion (when the possibility that the expansion has a tail end that consists entirely of the digit 9 is excluded). Therefore, the real number $r$ is not equal to any of $r_1, r_2, \ldots$ because the decimal expansion of $r$ differs from the decimal expansion of $r_i$ in the $i$th place to the right of the decimal point, for each $i$.

Because there is a real number $r$ between 0 and 1 that is not in the list, the assumption that all the real numbers between 0 and 1 could be listed must be false. Therefore, all the real numbers between 0 and 1 cannot be listed, so the set of real numbers between 0 and 1 is uncountable. Any set with an uncountable subset is uncountable. Hence, the set of real numbers is uncountable.
\end{thm}
\begin{thm}[uncountable supersets]
If $A$ and $B$ are countable sets, then $A \cup B$ is also countable.

\end{thm}

\begin{thm}[Schröder-Bernstein Theorem:] If $A$ and $B$ are sets with $|A| \leq |B|$ and $|B| \leq |A|$, then $|A| = |B|$. In other words, if there are one-to-one functions $f$ from $A$ to $B$ and $g$ from $B$ to $A$, then there is a one-to-one correspondence between $A$ and $B$.
\end{thm}
\begin{defn}[Computability]
We say that a function is computable if there is a computer program in some programming language that finds the values of this function. If a function is not computable, we say it is uncomputable.

\end{defn}
\newpage
\subsubsection{EXAMPLES}
\newpage

\section{Algorithms}
In class
\\
Please plan on signing up for this!
\newpage

\section{Number Theory and Cryptography}
\subsection{Divisibility and Modular Arithmetic}
\begin{defn}[Divides]
If \(a\) and \(b\) are integers with \(a \neq 0\), we say that \(a\) divides \(b\) if there is an integer \(c\) such that \(b = ac\), or equivalently, if \(\frac{b}{a}\) is an integer. When \(a\) divides \(b\), we say that \(a\) is a factor or divisor of \(b\), and that \(b\) is a multiple of \(a\). The notation \(a \mid b\) denotes that \(a\) divides \(b\). We write \(a \nmid b\) when \(a\) does not divide \(b\).

\end{defn}
\begin{thm}[basic properties of divisibility]
Let \(a\), \(b\), and \(c\) be integers, where \(a \neq 0\). Then
\begin{enumerate}
    \item[(i)] If \(a \mid b\) and \(a \mid c\), then \(a \mid (b + c)\);
    \item[(ii)] If \(a \mid b\), then \(a \mid bc\) for all integers \(c\);
    \item[(iii)] If \(a \mid b\) and \(b \mid c\), then \(a \mid c\).
\end{enumerate}

\end{thm}
\begin{thm}[useful consequence]
If \(a\), \(b\), and \(c\) are integers, where \(a \neq 0\), such that \(a \mid b\) and \(a \mid c\), then \(a \mid mb + nc\) whenever \(m\) and \(n\) are integers.


\end{thm}

\begin{thm}[The Division Algorithm]
Let \(a\) be an integer and \(d\) a positive integer. Then there are unique integers \(q\) and \(r\), with \(0 \leq r < d\), such that \(a = dq + r\).

\end{thm}
\begin{defn}[quotient, remainder]
In the equality given in the division algorithm, \(d\) is called the divisor, \(a\) is called the dividend, \(q\) is called the quotient, and \(r\) is called the remainder. This notation is used to express the quotient and remainder:
\[q = \frac{a}{d}, \quad r = a \mod d.\]


\end{defn}

\begin{defn}[modulo]
If \(a\) and \(b\) are integers and \(m\) is a positive integer, then \(a\) is congruent to \(b\) modulo \(m\) if \(m\) divides \(a - b\). We use the notation \(a \equiv b \pmod{m}\) to indicate that \(a\) is congruent to \(b\) modulo \(m\). We say that \(a \equiv b \pmod{m}\) is a congruence and that \(m\) is its modulus (plural moduli). If \(a\) and \(b\) are not congruent modulo \(m\), we write \(a \not\equiv b \pmod{m}\).


\end{defn}
\begin{thm}[modular equivalence]
Let \(a\) and \(b\) be integers, and let \(m\) be a positive integer. Then \(a \equiv b \pmod{m}\) if and only if \(a \mod m = b \mod m\).

\end{thm}
\begin{thm}[converting between modulo and mod]
Let \(m\) be a positive integer. The integers \(a\) and \(b\) are congruent modulo \(m\) if and only if there is an integer \(k\) such that \(a = b + km\).


\end{thm}
\begin{thm}[more conversions from mod to modulo]
Let \(m\) be a positive integer. If \(a \equiv b \pmod{m}\) and \(c \equiv d \pmod{m}\), then
\[a + c \equiv b + d \pmod{m}\] and \[ac \equiv bd \pmod{m}.\]

\end{thm}
\begin{thm}[a careful point]
Let \(m\) be a positive integer, and let \(a\) and \(b\) be integers. Then
\[(a + b) \mod m = ((a \mod m) + (b \mod m)) \mod m\]
and
\[ab \mod m = ((a \mod m)(b \mod m)) \mod m.\]

\end{thm}
\newpage
\subsubsection{EXAMPLES}
\begin{exmp}
Show that if \(a, b, c, \) and \(d\) are integers, where \(a \neq 0\), such that \(a \mid c\) and \(b \mid d\), then \(ab \mid cd\).
\begin{proof}$ $\\

\vspace{2cm}
\end{proof}
\end{exmp}

\begin{exmp}

Show that if \(a, b,\) and \(c\) are integers, where \(a \neq 0\) and \(c \neq 0\), such that \(ac \mid bc\), then \(a \mid b\).

\begin{proof}$ $\\

\vspace{2cm}
\end{proof}
\end{exmp}
\begin{exmp}

Prove or disprove that if \(a \mid bc\), where \(a, b,\) and \(c\) are positive integers and \(a \neq 0\), then \(a \mid b\) or \(a \mid c\).


\begin{proof}$ $\\

\vspace{2cm}
\end{proof}
\end{exmp}
\begin{exmp}
\begin{itemize} $ $\\
\item What are the quotient and remainder when 44 is divided by 8?
\item What time does a 12-hour clock read 80 hours after it reads 11:00?
\item Suppose that \(a\) and \(b\) are integers, \(a \equiv 4 \pmod{13}\), and \(b \equiv 9 \pmod{13}\). Find the integer \(c\) with \(0 \leq c \leq 12\) such that
\(c \equiv 2a + 3b \pmod{13}\)
\end{itemize}

\end{exmp}
\newpage
\subsection{Integer Representations and Algorithms}
\begin{thm}[Integer Representation]
Let \(b\) be an integer greater than 1. Then if \(n\) is a positive integer, it can be expressed uniquely in the form
\[n = a_k b^k + a_{k-1} b^{k-1} + \ldots + a_1 b + a_0,\]
where \(k\) is a nonnegative integer, \(a_0, a_1, \ldots, a_k\) are nonnegative integers less than \(b\), and \(a_k \neq 0\).
\end{thm}
\begin{defn}
The representation of \(n\) given in the Theorem above is called the base \(b\) expansion of \(n\). The base \(b\) expansion of \(n\) is denoted by \((a_k a_{k-1} \ldots a_1a_0)_b\). For instance, \((245)_8\) represents \(2 \cdot 8^2 + 4 \cdot 8 + 5 = 165\). 

\textbf{DECIMAL EXPANSIONS} Typically, the subscript 10 is omitted for base 10 expansions of integers because base 10, or decimal expansions, are commonly used to represent integers.

\textbf{BINARY EXPANSIONS:} Choosing 2 as the base gives binary expansions of integers. In binary notation, each digit is either a 0 or a 1. In other words, the binary expansion of an integer is just a bit string. Binary expansions (and related expansions that are variants of binary expansions) are used by computers to represent and do arithmetic with integers.

\textbf{OCTAL AND HEXADECIMAL EXPANSIONS:} Base 8 expansions are called octal expansions, and base 16 expansions are hexadecimal expansions.
The hexadecimal expansion of a number is a representation using the digits 0-9 and the letters A-F, where each digit contributes to the sum based on powers of 16. For example, $1A3_{16}$ corresponds to $1 \times 16^2 + 10 \times 16^1 + 3 \times 16^0$.

\end{defn}

\begin{exmp} $ $\\

\begin{itemize}
\item Convert the decimal expansion of each of these integers to a binary expansion. 231

\vspace{2cm}

\item  Convert the binary expansion of each of these integers to a decimal expansion.
 $(1 1111)_2$
 \vspace{2cm}

 \item Convert the hexadecimal expansion of each of these integers
to a binary expansion. $(80E)_{16}$

\newpage
\item Convert $(100110111)_2$ to octal.
\newline
\[\sum_{i=0}^8 a_i*2^i\approx \sum_{i=0}^3 b_i*(2^3)^i \]

\begin{align}
=&(a_8a_7a_6a_5a_4a_3a_2a_1a_0)_2\\
=&\sum_{i=0}^8a_i*2^i\\
=&(a_84+a_72+a_6)(2^3)^2+(a_54+a_42+a_3)2^3+a_2*4+a_1*2+a_0\\
=&(b_2)(2^3)^2+(b_1)2^3+b_0\\
=&\sum_{i=0}^2 b_i*(2^3)^i\\
=&(b_2b_1b_0)_{8}
\end{align}

\[(110110111)_2=(667)_{8}\]


\end{itemize}

\end{exmp}

\newpage
\subsection{Primes and Greatest Divisors}
\begin{defn}[prime]
An integer $p$ greater than 1 is called prime if the only positive factors of $p$ are 1 and $p$.
A positive integer that is greater than 1 and is not prime is called composite.
\end{defn}

\begin{thm}{THE FUNDAMENTAL THEOREM OF ARITHMETIC}$ $\\
Every integer greater than 1 can be written uniquely as a prime or as the product of two or more primes, where the prime factors are written in order of nondecreasing size.
\end{thm}
\begin{thm}[Factor Bounds]$ $\\
If $n$ is a composite integer, then $n$ has a prime divisor less than or equal to $\sqrt{n}$.

\end{thm}
\begin{exmp}Show that 101 is prime.
\begin{proof} $ $\\
\vspace{2cm}
 $ $\\
 101 is not divisible by 2, 3, 5, 7, 11. By Theorem 4.20, since 101 is not divisible by primes less that root 101, we are good!
\end{proof}
\end{exmp}
\begin{exmp}{The Sieve of Eratosthenes}$ $\\

The Sieve of Eratosthenes is like an old-school trick for figuring out all the prime numbers up to a certain point. You start with 2 and cross out all its multiples, then move to the next unmarked number and do the same. Keep going until you've dealt with the multiples up to the square root of the limit – what's left are the primes.
\begin{enumerate}
  \item \textbf{Initialization:}
    \begin{itemize}
      \item Create a list of boolean values representing numbers from 2 to the given limit, initially all marked as true.
      \item Mark 0 and 1 as false since they are not prime.
    \end{itemize}
  
  \item \textbf{Iteration:}
    \begin{itemize}
      \item Start with the first number in the list (2), the smallest prime.
      \item Mark all multiples of the current prime as false, as they cannot be prime.
      \item Move to the next number in the list that is still marked as true, and repeat the process.
      \item Continue until the square of the current prime exceeds the given limit.
    \end{itemize}
  
  \item \textbf{Completion:}
    \begin{itemize}
      \item The remaining unmarked numbers in the list are primes.
    \end{itemize}
  
  \item \textbf{Result:}
    \begin{itemize}
      \item The algorithm produces a list of all prime numbers up to the given limit.
    \end{itemize}
\end{enumerate}
\href{https://pythontutor.com/render.html#code=def%20sieve_of_eratosthenes%28limit%29%3A%0A%20%20%20%20primes%20%3D%20%5B%5D%0A%20%20%20%20is_prime%20%3D%20%5B1%5D%20*%20%28limit%20%2B%201%29%0A%20%20%20%20is_prime%5B0%5D%20%3D%20is_prime%5B1%5D%20%3D%200%0A%0A%20%20%20%20for%20number%20in%20range%282,%20int%28limit**0.5%29%20%2B%201%29%3A%0A%20%20%20%20%20%20%20%20if%20is_prime%5Bnumber%5D%3A%0A%20%20%20%20%20%20%20%20%20%20%20%20primes.append%28number%29%0A%20%20%20%20%20%20%20%20%20%20%20%20for%20multiple%20in%20range%28number%20*%20number,%20limit%20%2B%201,%20number%29%3A%0A%20%20%20%20%20%20%20%20%20%20%20%20%20%20%20%20is_prime%5Bmultiple%5D%20%3D%200%0A%0A%20%20%20%20for%20number%20in%20range%28int%28limit**0.5%29%20%2B%201,%20limit%20%2B%201%29%3A%0A%20%20%20%20%20%20%20%20if%20is_prime%5Bnumber%5D%3A%0A%20%20%20%20%20%20%20%20%20%20%20%20primes.append%28number%29%0A%0A%20%20%20%20return%20primes%0A%0A%23%20Example%20usage%3A%0Alimit%20%3D%2020%0Aresult%20%3D%20sieve_of_eratosthenes%28limit%29%0Aprint%28%22Prime%20numbers%20up%20to%22,%20limit,%20%22are%3A%22,%20result%29&cumulative=true&curInstr=0&heapPrimitives=false&mode=display&origin=opt-frontend.js&py=3&rawInputLstJSON=%5B%5D&textReferences=true}{Example implementation here}
\end{exmp}

\begin{thm}{Infinitude of primes} 
There are infinitely many primes.
\begin{proof}
$ $\\
\vspace{1cm}
Suppose that there are a finite number of primes $p_0,p_1,p_2,p_3,\cdots,p_n$.
\newline
$n=p_0\cdot p_1\cdot p_2 \cdot p_3 \cdot\cdots\cdot p_n$
\newline
$n\% p_i=0$
\newline
$n+1\% p_i=1$
\newline
$n+1$ is a new prime! this contradicts our assumption, therefore there are an infinite number of primes.
$ $\\
\end{proof}
\end{thm}
\begin{defn}{GCD}$ $\\
Let $a$ and $b$ be integers, not both zero. The largest integer $d$ such that $d | a$ and $d | b$ is called the greatest common divisor of $a$ and $b$. The greatest common divisor of $a$ and $b$ is denoted by $\gcd(a, b)$.
\end{defn}
\begin{exmp}Find $\gcd(48, 18)$
$ $\\
Since 6 is the largest divisor of $18$ which divides $48$ 6 is the $\gcd(48,18)$
\vspace{0.5em}
\noindent
$18=2*3*3$\newline
$48=2*2*2*2*3$
\end{exmp}


\begin{defn}{Relatively Prime}$ $\\
The integers $a$ and $b$ are relatively prime if their greatest common divisor is 1.
\end{defn}
\begin{exmp} Find which pairs of the following numbers are relatively prime. 
25,9,12, and 7.
\newline
\noindent
$25=5*5$\newline
$9=3*3$\newline
$12=3*2*2$\newline
$7=7$\newline


\end{exmp}
\begin{defn}{LCM}$ $\\
The least common multiple of the positive integers $a$ and $b$ is the smallest positive integer that is divisible by both $a$ and $b$. The least common multiple of $a$ and $b$ is denoted by $\lcm(a, b)$.
\end{defn}
\begin{exmp}
Compute $\lcm(6,9)$
\newline
$6=2*3$
\newline
$9=3*3$

\end{exmp}
\begin{thm}{LCM and GCD}$ $\\
Let $a$ and $b$ be positive integers. Then $ab = \gcd(a, b) \cdot \lcm(a, b)$.
\end{thm}
\newpage
\begin{exmp}
Find $\gcd(12,15)$ and $\gcd(12,15)$ and  confirm $12\cdot15 = \gcd(12, 15) \cdot \lcm(12, 15)$.
\newline
$12=3*2*2$\newline
$15=5*3$\newline
$\gcd(12,15)=3$\newline
$\lcm(12,15)=5*3*2*2$\newline
\newline
$\lcm(a,b)=\frac{a\cdot b}{\gcd(a,b)}$


\end{exmp}
\begin{exmp}Suppose $a|b$, what is $\gcd(a,b)$? what is $\lcm(a,b)$
$ $\\
$a=\gcd(a,b)$\\

\end{exmp}
\
\begin{thm}{Lemma for Euclidean Algorithm}$ $\\
Let $a = bq + r$, where $a$, $b$, $q$, and $r$ are integers. Then $\gcd(a, b) = \gcd(b, r)$.

\end{thm}
\newpage
\begin{thm}{Euclidean Algorithm}$ $\\
The Euclidean Algorithm is a method for finding the greatest common divisor (gcd) of two integers. Let $a$ and $b$ be positive integers. The algorithm proceeds by iteratively replacing the larger number with the remainder of the division of the larger number by the smaller number. This process continues until the remainder is zero. The gcd is then the last non-zero remainder. The algorithm can be expressed as follows:

\textbf{Euclidean Algorithm:}

\begin{enumerate}
  \item \textbf{Initialization:}
    \begin{itemize}
      \item Let $a$ and $b$ be positive integers.
    \end{itemize}
  
  \item \textbf{Iteration:}
    \begin{itemize}
      \item Replace the larger number $a$ with the remainder of $a$ divided by $b$.
      \item Continue this process until the remainder is zero.
      \item The greatest common divisor (gcd) is the last non-zero remainder, which is $b$.
    \end{itemize}
  
  \item \textbf{Result:}
    \begin{itemize}
      \item The algorithm efficiently computes the gcd of $a$ and $b$.
    \end{itemize}
\end{enumerate}

\end{thm}
\begin{exmp}Use the Euclidean algorithm to show 48 and 
18. Write down all of the quotients and remainders which were calculated.
\vspace{3cm}
\end{exmp}
\href{https://pythontutor.com/render.html#code=def%20euclidean_algorithm%28a,%20b%29%3A%0A%20%20%20%20while%20b%20!%3D%200%3A%0A%20%20%20%20%20%20%20%20print%28f%22%7Ba%7D%3D%7Bb%7Dx%7Ba//b%7D%2B%7Ba%25b%7D%22%29%0A%20%20%20%20%20%20%20%20a,%20b%20%3D%20b,%20a%20%25%20b%0A%20%20%20%20return%20a%0A%0A%23%20Example%20usage%3A%0Anum1%20%3D%2048%0Anum2%20%3D%2018%0Aresult%20%3D%20euclidean_algorithm%28num1,%20num2%29%0A%0Aprint%28f%22The%20gcd%20of%20%7Bnum1%7D%20and%20%7Bnum2%7D%20is%3A%20%7Bresult%7D%22%29&cumulative=true&curInstr=18&heapPrimitives=nevernest&mode=display&origin=opt-frontend.js&py=3&rawInputLstJSON=%5B%5D&textReferences=false}{Implementation Here}
\begin{thm}{Bézout's Theorem}$ $\\
If $a$ and $b$ are positive integers, then there exist integers $s$ and $t$ such that $\gcd(a, b) = sa + tb$.
\end{thm}
\begin{exmp}
Find $s$ and $t$ such that $\gcd(48,18)= s48 + t18$
\vspace{2cm}
\newline
48=18x2+12
\newline
18=12x1+6
\newline
12=6x2+0
\newline
6=18x3-48
\end{exmp}

\begin{thm}{Prime Factors}$ $\\
If $p$ is a prime and $p \mid a_1 \cdot a_2 \cdot \ldots \cdot a_n$, where each $a_i$ is an integer, then $p \mid a_i$ for some $i$.
\end{thm}
\begin{exmp}{Find the prime factorization of 360}
\vspace{2cm} $ $\\
360=2*2*2*3*3*5
\end{exmp}

\begin{exmp}
$3+(-3)=0$ additive ($a+0=a$)!
\newline
$2*(.5)=1$ multiplicative ($a*1=a$)
\newline
$(6*1)\mod p$. $1$ is multiplicative identity when dealing with mods.
\newline
$(2*s)=1 (\mod 3)$ 2 is multiplicative inverse of itself.
\newline
$(2*s)=1 (\mod 5)$ 3 is the multiplicative inverse!
\newline
$(2*s)=1 (\mod 7)$ 4 is the multiplicative inverse
\newline
$(2*s)=1 (\mod 6)$ DOES NOT EXIST!!!!!.
\newline

\end{exmp}
\begin{thm}{Multiplicative Inverse Property}

If $a$ and $m$ are relatively prime integers, where $m > 1$, then an inverse of $a$ modulo $m$ exists. Furthermore, this inverse is unique modulo $m$. In other words, there is a unique positive integer $a^{-1}$ less than $m$ that is an inverse of $a$ modulo $m$, and every other inverse of $a$ modulo $m$ is congruent to $a^{-1}$ modulo $m$.

\end{thm}
\begin{exmp} Does the modular inverse of $22 \mod 5$ exist? explain why. Then, use Bézout's Theorem to find $s$ and $t$ such that $s\cdot 22+t\cdot 5=1$, take both sides $\mod 5$ and find the inverse of $22 \mod 5$.
\newline
$s\cdot22+t\cdot 5=1$
\newline
22-5*4=2
\newline
5-2x2=1
\newline
5-(22-5*4)x2=1
\newline $ $\\
5*9-22*2=1
\newline
$(5*9-22*2)\mod 5=1\mod 5$
\newline
$(22*-2)\mod 5=1\mod 5$
\newline
modular inverse is $3$
\newline


\end{exmp}



\newpage
\subsection{Solving Congruences}$ $\\
\begin{thm}{Chinese Remainder Theorem:}


Let $m_1, m_2, \ldots, m_k$ be pairwise coprime positive integers, and let $a_1, a_2, \ldots, a_k$ be integers. Then, the system of congruences:


\begin{align*}
x &\equiv a_1 \pmod{m_1} \\
x &\equiv a_2 \pmod{m_2} \\
& \vdots \\
x &\equiv a_k \pmod{m_k}
\end{align*}


has a unique solution modulo $M = m_1 \cdot m_2 \cdot \ldots \cdot m_k$ given by:

\[
x \equiv a_1 M_1 y_1 + a_2 M_2 y_2 + \ldots + a_k M_k y_k \pmod{M},
\]

where $M_i = M/m_i$ and $y_i$ is the modular inverse of $M_i$ modulo $m_i$.

\end{thm}
\begin{exmp}
Find a value of $x$ which satisfies the system of congruences:

\begin{align*}
x &\equiv 2 \pmod{3} \\
x &\equiv 3 \pmod{5} \\
x &\equiv 1 \pmod{7}
\end{align*}


The Chinese Remainder Theorem allows us to find a unique solution modulo \(3 \times 5 \times 7 = 105\).

Let \(M_1 = 5 \times 7 = 35\), \(M_2 = 3 \times 7 = 21\), and \(M_3 = 3 \times 5 = 15\).

Find the modular inverses:

\begin{align*}
35 \cdot y_1 &\equiv 1 \pmod{3} \quad \text{(solve for \(y_1\))} \quad \Rightarrow \quad y_1 \equiv 2 \pmod{3} \\
21 \cdot y_2 &\equiv 1 \pmod{5} \quad \text{(solve for \(y_2\))} \quad \Rightarrow \quad y_2 \equiv 1 \pmod{5} \\
15 \cdot y_3 &\equiv 1 \pmod{7} \quad \text{(solve for \(y_3\))} \quad \Rightarrow \quad y_3 \equiv 1 \pmod{7}
\end{align*}


Apply the Chinese Remainder Theorem formula:
\[ x \equiv (2 \cdot 35 \cdot 2) + (3 \cdot 21 \cdot 1) + (1 \cdot 15 \cdot 1) \pmod{105} \]

Solving this congruence, we get \(x \equiv 41 \pmod{105}\).
\end{exmp}
\begin{thm}{Fermat's Little Theorem}$ $\\

If $p$ is prime and $a$ is an integer not divisible by $p$, then $a^{p-1} \equiv 1 \pmod{p}$. Furthermore, for every integer $a$, we have $a^p \equiv a \pmod{p}$.

\end{thm}
\newpage


\section{Induction} $ $\\

Is a tomato a fruit? Perhaps.
\newline
Who decides what isn't a fruit and what is? Botanists may say that tomatoes are a fruit because they from a flower and contain seeds. Botanists don't care how tomatoes taste in a fruit salad. 
\newline
When we are defining combinatorial structures and talking about them, we care about aspects of these objects that perhaps others using them are not particularly interested in. Botanists define fruit in a way that make it very clear what is and isn't a fruit (according to them anyway).
\newline
The following conversation regarding the construction of the positive natural numbers is a bit like botanists discussing "what is a fruit and how do we define it".
\newline

Virtually all of the infinite sets that we have constructed so far have been based off the existence of a well understood infinite set, the natural numbers. We have not yet formally defined or constructed them. Is there only one set of Natural numbers? Are there many? Is it up to interpretations? Below are a handful of observations that led to the construction we all use today.
\newline

We have a pretty good idea of what the positive natural numbers are $1,2,\cdots $, but what the hell does the $\cdots$ mean? Is infinity a natural number then? $\infty\in \mathbb{N}$? Are there an infinite number of natural numbers?
\newline
We do have an idea of a next natural number, we define it as the successor.
\begin{defn}{successor}
If $n$ is a positive integer, we call $S(n)$ the successor of $n$ and we say $n<S(n)$. We write $S(n)=n+1$.
\end{defn}
We also have an idea of what it means for one natural number to be bigger than another. 
\begin{defn}{$<$} $ $ \\
We define the binary operation $<:\mathfrak{U}\times \mathfrak{U}\to \{\text{True},\text{False}\}$ where $ a < S(a)$ and $<$ satisfies the transitive property $a<b\land b<c \to a <c$. If $a < b$ we say $a$ is less than $b$.
\end{defn}
Now we can try to define the natural numbers:
\begin{exmp}[Constructing infinite sets]
Consider a set $N$ with the following properties:
\begin{enumerate}
\item It contains $1$
\item If $n$ is in the set, then the successor of $n$ is a natural number.
\item All positive integers are the successor of a unique positive integer or are $1$.
\end{enumerate}
We can imagine the members of these sets as nodes, or dots. We can draw an arrow from a dot $a$ to a dot $b$ if the successor of $a$ is $b$.
\newline
Do we all have the same drawing? Is a set defined this way unique? or can our drawings look different? Can you prove to me that there is no Natural number that ends in a $.7$? It doesn't seem like there should be one?
\newline
\vspace{2cm}
\end{exmp}


\newpage
As it turns out, this is not enough. We can have all sorts of weird looking collections of dots that don't break any of these rules. We need one additional rule:

\begin{defn}[The well ordering principle]
Every nonempty subset of the set of positive
integers has a least element.
\end{defn}
\begin{exmp} Prove that no natural number ends in a $.7$.
\vspace{3cm}
\end{exmp}


As it turns out, the well ordering principle is equivalent to the axiom of mathematical induction.

\begin{defn}[Axiom of Mathematical Induction]
If $S$ is a set of positive integers such that $1 \in S$ and
for all positive integers $n$ if $n \in S$, then $n + 1 \in S$, then $S$ is the set of positive integers.
\end{defn}
We can use mathematical induction to prove the above example as well.
\begin{exmp} Prove that no natural number ends in a $.7$.
\vspace{3cm}
\end{exmp}



\begin{defn}[Induction] $ $\\
Inference Rule:
\newline
For any proposition $P$
\newline
If $P(0)$, $\forall n \in \mathbb{N},P(n) \to P(n+1)$, then $\forall n \in \mathbb{N}, P(n)$
\newline
We say $P(0)$ is a base case.
\newline
Inductive Step: $P(n) \to P(n+1)$
\newline
Inductive Hypothesis: $P(n)$
\end{defn}

\begin{exmp}
Use induction to show that $\forall n\in \mathbb{N}, n< 2^n$.
\begin{proof}
Let $P(n)=:\text{" $ n< 2^n$"}$
\newline
We proceed by induction:
\newline
Base Case, $n=0$
\newline
Since $0<1$ we see $P(0)$ is true. 
\newline
Inductive Step:
\newline
Assume the Inductive Hypothesis (I.H.)
\newline
Goal:
\newline
We assume $P(n)=\text{"$n<2^n$"}$ is true
\newline
W.T.S. $P(n+1)=:\text{" $ (n+1)< 2^{(n+1)}$"}$
\newline
\begin{align}
(n+1)&=n+1\\
    &<2^n+1 \text{ By I.H.}\\
    &=2^n+2^0\\
    &\leq 2^n+2^n
    &\leq 2*(2^n)\\
    &=2^{n+1}
\end{align}
By induction we see that,
$\forall n \in \mathbb{N}, n<2^n$

\end{proof}

\end{exmp}
\begin{exmp}
Use induction to show that $n^3-n$ is divisible by 3 whenever $n$ is a positive integer. 
\begin{proof}
Let $P(n)=\text{"$n^3-n$ is divisible by 3"}$
\newline

\end{proof}
We proceed by induction
\newline
Base Case:
\newline
We see that $P(1)$ is true since 0 is divisible by 3.
\newline
Inductive Step: $P(n)\to P(n+1)$
\newline
Inductive Hypothesis: $P(n)$
\newline
Assume $n^3-n$ is divisible by 3.
\newline
W.T.S. $P(n+1)$
\newline
$(n+1)^3-(n+1)$ is divisible by 3.
\newline
\begin{align}
(n+1)^3-(n+1)=&(n+1)(n+1)(n+1)-(n+1)\\
=&n^3+3n^2+3n+1-(n+1)\\
=&n^3-n +3n^2+3n+1-(1)\\
=&n^3-n +3n^2+3n\\
=&n^3-n +3(n^2+n)
\end{align}
By Inductive Hypothesis $n^3-n$, and since $3(n^2+n)$ is divisible by three, their sum is divisible by 3 and we have show $P(n+1)$.
\end{exmp}
\newpage

\section{Counting}
\subsection{The Basics of Counting}
 In the back of your mind you will want to be thinking of this as considering Sets and their cardinality.

\begin{defn}[Product Rule]
    Suppose that a procedure can be broken down into a sequence of two tasks. If there are $n_1$ ways to do the first task and for each of these ways of doing the first task, there are $n_2$ ways to do the second task, then there are $n_1 n_2$ ways to do the procedure.
\end{defn}
\begin{exmp}[Set analog] Recall that $|A\times B|=|A|\times |B|$
\end{exmp}
\begin{exmp}
    A new company, XYZ Corp, has just two employees, Sanchez and Patel. They have rented a floor of a building with 12 offices. The company is deciding how to assign these offices to the employees. How many ways are there to assign different offices to Sanchez and Patel?
    
    \vspace{2cm}
\end{exmp}
\begin{exmp}
    The chairs of an auditorium are to be labeled with an uppercase English letter followed by a positive integer not exceeding 100. What is the largest number of chairs that can be labeled differently?
    \vspace{2cm}
\end{exmp}
\begin{exmp}
    How many different bit strings of length seven are there?
    \vspace{2cm}
\end{exmp}
\begin{exmp}
    Counting Functions: How many functions are there from a set with $m$ elements to a set with $n$ elements?
    \vspace{2cm}
\end{exmp}
\begin{defn}[Sum Rule]
    If a task can be done either in one of $n_1$ ways or in one of $n_2$ ways, where none of the set of $n_1$ ways is the same as any of the set of $n_2$ ways, then there are $n_1 + n_2$ ways to do the task.
\end{defn}
\begin{exmp}[Set analog] Recall that if $A$ and $B$ are disjoint sets,  $|A\cup B|=|A|+|B|$
\end{exmp}

\begin{exmp}
    Suppose that either a member of the CS faculty or a grad student is chosen as a representative to a university committee. How many different choices are there for this representative if there are 37 members of the CS faculty and 83 grad students, and no one is both a faculty member and a student?

    \vspace{2cm}
\end{exmp}

\newpage
\begin{exmp}
    What is the value of $k$ after the following code, where $n_1, n_2, \ldots, n_m$ are positive integers, has been executed?

    \begin{verbatim}
k := 0
for i_1 := 1 to n_1
    k := k + 1
for i_2 := 1 to n_2
    k := k + 1
.
.
.
for i_m := 1 to n_m
    k := k + 1
\end{verbatim}
\vspace{2cm}
\end{exmp}
\subsubsection{Tricky Counting Problems}
\begin{exmp}
    Each user on a computer system has a password, which is six to eight characters long, where each character is an uppercase letter or a digit. Each password must contain at least one digit. How many possible passwords are there?
    \vspace{4cm}
\end{exmp}

\begin{defn}[Subtraction Rule]
    If a task can be done in either $n_1$ ways or $n_2$ ways, then the number of ways to do the task is $n_1 + n_2$ minus the number of ways to do the task that are common to the two different ways.
\end{defn}
\begin{exmp}[set analog]
For any sets $A$ and $B$,
$|A\cup B|=|A|+|B|-|A\cap B|$ (this is called the principle of inclusion-exclusion)
\end{exmp}
\begin{exmp}
If I have 6 shirts and 3 are red, and I have 3 hats and 2 are  red, how many ways can I get dressed and make sure that I am wearing red?
\vspace{2cm}
\end{exmp}

\begin{defn}[Division Rule]
    There are $\frac{n}{d}$ ways to do a task if it can be done using a procedure that can be carried out in $n$ ways, and for every way $w$, exactly $d$ of the $n$ ways correspond to way $w$.
\end{defn}
\begin{exmp}[Set Analog]
If you take a set of size $n$ and break this down into a collection of disjoint sets of size $d$, then there are $\frac{n}{d}$ of these sets.

\end{exmp}
\begin{exmp}
    How many different ways are there to seat four people around a circular table, where two seatings are considered the same when each person has the same left neighbor and the same right neighbor?

    \vspace{2cm}
\end{exmp}
\subsection{The Pigeonhole Principle}

\begin{defn}[Pigeonhole Principle]
    If $k$ is a positive integer and $k + 1$ or more objects are placed into $k$ boxes, then there is at least one box containing two or more of the objects.
\end{defn}
\begin{exmp}

    A function $f$ from a set with $k + 1$ or more elements to a set with $k$ elements is not one-to-one.

\end{exmp}
\begin{thm}[Generalized Pigeonhole Principle]
    If $N$ objects are placed into $k$ boxes, then there is at least one box containing at least $\lceil N/k \rceil$ objects.
\end{thm}
\begin{proof} $ $\\
\vspace{4cm}
$ $\\
\end{proof}

\begin{exmp} Suppose that between 10 people, they have 31 marbles. Prove that at least one person has at least 4 marbles.
\vspace{2cm}
\end{exmp}

\begin{exmp} Suppose that on average, some group of people has 3.1 marbles. Prove that at least one person has at least 4 marbles.
\vspace{2cm}
\end{exmp}
\begin{exmp}
    How many cards must be selected from a standard deck of 52 cards to guarantee that at least three cards of the same suit are chosen?
\vspace{3cm}
\end{exmp}
\subsubsection{Ramsey Problems}
\begin{exmp}
    Assume that in a group of six people, each pair of individuals consists of two friends or two enemies. Show that there are either three mutual friends or three mutual enemies in the group.
    \vspace{3cm}
\end{exmp}

\newpage

\subsection{Permutations and Combinations}
\begin{defn}[Permutation]
A permutation of a set of distinct objects is an ordered arrangement of these objects.
We also are interested in ordered arrangements of some of the elements of a set. An ordered
arrangement of r elements of a set is called an r-permutation.
    The number of $r$-permutations of a set with $n$ elements is denoted by $P(n, r)$.


\end{defn}
\begin{exmp}
    Let $S = \{1, 2, 3\}$. The ordered arrangement $3, 1, 2$ is a permutation of $S$. The ordered arrangement $3, 2$ is a $2$-permutation of $S$.
\end{exmp}

\begin{thm}
    If $n$ is a positive integer and $r$ is an integer with $1 \leq r \leq n$, then there are $P(n, r) = n(n - 1)(n - 2) \cdot \ldots \cdot (n - r + 1)$ $r$-permutations of a set with $n$ distinct elements.
\end{thm}
We will almost always use this notation:
\begin{defn}
    If $n$ and $r$ are integers with $0 \leq r \leq n$, then $P(n, r) = \frac{n!}{(n - r)!}$. We call this "$n$ pick $r$".
\end{defn}
\begin{exmp}
How many different ways can we select a first-prize winner, a second-prize winner, and a third-prize winner from a pool of 100 contestants who have entered a contest?
\vspace{2cm}
\end{exmp}


\begin{exmp}

How many permutations of the letters ABCDEFGH contain the string ABC.
\vspace{3cm}
\end{exmp}
\begin{defn}
    An $r$-combination of elements of a set is an unordered selection of $r$ elements from the set. Thus, an $r$-combination is simply a subset of the set with $r$ elements.
\end{defn}
\begin{defn}
    The number of $r$-combinations of a set with $n$ distinct elements is denoted by $C(n, r)$. Note that $C(n, r)$ is also denoted by $\binom{n}{r}$ and is called a binomial coefficient.
\end{defn}
\begin{thm}
    The number of $r$-combinations of a set with $n$ elements, where $n$ is a nonnegative integer and $r$ is an integer with $0 \leq r \leq n$, equals
    \[
    C(n, r) = \frac{n!}{r! (n - r)!}.
    \]
\end{thm}
\begin{exmp}
    The number of permutations $P(n, r)$ is equal to the product of the number of combinations $C(n, r)$ and the number of permutations $P(r, r)$:
    \[
    P(n, r) = C(n, r) \cdot P(r, r).
    \]
    \vspace{3cm}
\end{exmp}
\begin{thm}

    Let $n$ and $r$ be nonnegative integers with $r \leq n$. Then, 
    \[
    C(n, r) = C(n, n - r).
    \]
    \vspace{3cm}
\end{thm}
\subsection{Binomial Coefficients}

First Consider Pascal's Identity:
\begin{thm}[Pascal's Identity]
    Let \(n\) and \(k\) be positive integers with \(n \geq k\). Then
    \[
    \binom{n+1}{k} = \binom{n}{k-1} + \binom{n}{k}.
    \]
\begin{proof}
Pascals' Identity claims that the numbers of ways a committee of $k$ people can be formed from a collection of $n+1$ people is the same as the number of ways a committee of $k-1$ people can be formed from a group of $n$ people plus the numbers of ways a committee of $k$ people can be formed by a committee of $n$ people.
\end{proof}
\end{thm}

\begin{exmp}
For example:
Suppose that I have a set of size 5, $S$. I would like to count the number of subsets of size 3. This is $\binom{5}{3}$.
\newline
For $S=\{A,B,C,D,E\}$, we have 
$\{A,B,C\}$,$\{A,B,D\}, \{A,B,E\}$,
$\{A,C,D\}$,$\{A,C,E\}$,$\{A,D,E\}$,
$\{B,C,D\},\{B,C,E\},\{B,D,E\}$
\newline
so there are nine in total. 
\newline
What is another way of counting these 9?
\newline
I can first consider all the subsets of size 3 which contain a $A$: $\{A,B,C\}$,$\{A,B,D\}, \{A,B,E\}$,
$\{A,C,D\}$,$\{A,C,E\}$ and $\{A,D,E\}$
and then consider all the subsets which don't contain $A$: $\{B,C,D\},\{B,C,E\},\{B,D,E\}$.
\newline
The number of subsets of size 3 which include $A$ have two more choices of remaining values. Therefore these choices are equivalent to choosing two elements from $\{B,C,D,E\}$--$\binom{4}{2}$.
\newline
The number of subsets of size 3 which do not include $A$ have three choices of remaining values. Therefore these choices are equivalent to choosing three elements from $\{B,C,D,E\}$--$\binom{4}{3}$.
\end{exmp}
 From this observation we can create a Triangle called Pascals triangle. 


\begin{exmp}
We will be using the notation $\binom{n}{k}$ to represent $C(n,k)$ (the number of ways we can select $k$ unordered distinct things from a collection of $n$ distinct things)
For example, from above we have:
\newline
$\binom{3}{2}=\binom{2}{2}+\binom{2}{1}$
\newline
we can also see that $\binom{2}{1}=\binom{1}{1}+\binom{1}{0}$
\newline
Note that $\binom{n}{0}=1$ and $\binom{n}{n}=1$.
\newline
In this way we see a RECURSIVE relationship between binomial terms.\\
We can then write $\binom{n}{k}$ as a recursive function below:
\newline
\begin{verbatim}
def binom(n,k):
    if n==k:
        return 1
    elif k==0:
        return 1
    else:
        return binom(n-1,k)+binom(n-1,k-1)
print(binom(4,2))
\end{verbatim}
\end{exmp}

\begin{exmp}
This recursive relationship can also be seen in\\
Pascals Triangle. The first few rows are the values
\newline
$\binom{0}{0}$
\newline
$\binom{1}{0}$,$\binom{1}{1}$
\newline
$\binom{2}{0}$, $\binom{2}{1}$,$\binom{2}{2}$
\newline
Can you use Pascal's Triangle to calculate $\binom{6}{3}$?

\begin{verbatim}
              1
            1   1
          1   2   1
        1   3   3   1
      1   4   6   4   1
    1   5  10  10   5   1
  1   6  15  20  15   6   1
\end{verbatim}
\end{exmp}
As it turns out, this observation gives us the binomial theorem. Which we use to understand a probability distribution called the binomial distribution-- it is the discrete version of the normal distribution (perhaps the most well known probability distribution).
\begin{thm}[The Binomial Theorem]
    Let \(x\) and \(y\) be variables, and let \(n\) be a nonnegative integer. Then
    \[
    (x + y)^n = \sum_{j=0}^{n} \binom{n}{j} x^{n-j}y^j = \binom{n}{0}x^n + \binom{n}{1}x^{n-1}y + \ldots + \binom{n}{n-1}xy^{n-1} + \binom{n}{n}y^n.
    \]
\end{thm}

\subsubsection{Permutations with Repetition}
\begin{thm}
    The number of $r$-permutations of a set of $n$ objects with repetition allowed is $n^r$.
\end{thm}
\begin{exmp}
    How many strings of length $r$ can be formed from the uppercase letters of the English alphabet?
\vspace{3cm}
\end{exmp}
\subsubsection{Combinations with Repetition}
\begin{thm}
    There are $\binom{n + r - 1}{r} = \binom{n + r - 1}{n - 1}$ $r$-combinations from a set with $n$ elements when repetition of elements is allowed. (This is called the starts and bars argument).
\end{thm}
\begin{exmp}
    Suppose that a cookie shop has four different kinds of cookies. How many different ways can six cookies be chosen? Assume that only the type of cookie, and not the individual cookies or the order in which they are chosen, matters.
\vspace{3cm}
\end{exmp}
\subsubsection{Permutations with Indistinguishable Objects}

\begin{thm}
The number of different permutations of $n$ objects, where there are $n_1$ indistinguishable objects of type 1, $n_2$ indistinguishable objects of type 2, ..., and $n_k$ indistinguishable objects of type $k$, is
\[
\frac{n!}{n_1! \cdot n_2! \cdot \ldots \cdot n_k!}.
\]

\end{thm}

\begin{exmp}
How many different strings can be made by reordering the letters of the word SUCCESS?\vspace{3cm}
\end{exmp}


\begin{exmp}
How many ways are there to distribute hands of 5 cards to each of four players from the standard
deck of 52 cards?
\end{exmp}

\section{Discrete Probability}
\subsection{Intro to Discrete Probability}
\begin{defn}
    If $S$ is a finite nonempty sample space of equally likely outcomes, and $E$ is an event, i.e., a subset of $S$, then the probability of $E$ is given by $P(E) = \frac{|E|}{|S|}$.
\end{defn}

\begin{exmp}
If the sample space $S=\{a,b,c,d,e\}$ and the $E=:\{x| x \text{ is a vowel}\}$, then $P(E)=\frac{|\{a,e\}|}{|S|}=\frac{2}{5}$. 
\newline
This is one way that we model probability. We would say something like, if we randomly pick a letter from $S$, the probability that it is a vowel is $\frac{2}{5}$.
\newline

\end{exmp}

\begin{exmp}
    Let $E$ be an event in a sample space $S$. The probability of the event $E$'s complement, denoted as $S - E$, is given by
    \[
    P(S-E) = 1 - P(E).
    \]
\end{exmp}
\begin{exmp}
For instance, in the example above if we want to know the probability of randomly selecting a consonant, we can instead subtract the probability of finding a vowel from one!
\newline
Let $C=\{b,c,d\}$ then notice that $C=\overline{S}$ and $P(C)=\frac{3}{5}=1-\frac{2}{5}$.
\newline
\end{exmp}
\begin{exmp}
If the probability that you eat chicken tomorrow is $.3$ what is the probability you don't eat chicken tomorrow?
\vspace{2cm}
\end{exmp}

\begin{exmp}
    Let $E_1$ and $E_2$ be events in the sample space $S$. Then the probability of the union of events $E_1$ and $E_2$ is given by
    \[
    p(E_1 \cup E_2) = p(E_1) + p(E_2) - p(E_1 \cap E_2).
    \]
\end{exmp}
\begin{exmp}
If the probability you eat breakfast at your mom's house is .2 and the probability you eat breakfast at a diner is .3 and you are not a hobbit (you only eat one breakfast).
What is the probability that you eat breakfast at your mom's or at a diner?
\vspace{2cm}
\end{exmp}
\begin{exmp}[The Monty Hall Three-Door Puzzle] $ $\\ Suppose you are a game show contestant. You have a
chance to win a large prize. You are asked to select one of three doors to open; the new car
is behind one of the three doors. Behind another door is a goat. Behind the third door is trash. Once you select a door, the
game show host, who knows what is behind each door, does the following. First, whether or not
you selected the winning door (car), he opens one of the other two doors that he knows is a losing
door (selecting at random if both are losing doors). Then he asks you whether you would like to
switch doors. Which strategy should you use? Should you change doors or keep your original
selection, or does it not matter?

\noindent Let $G={\text{"you picked the goat door"}}$
\newline
Let $T={\text{"you picked the trash door"}}$
\newline
Let $C={\text{"you picked the car door"}}$
\newline
Draw a decision tree of the scenario that you keep your choice, and find the probability of winning the car and then compare this to the decision tree of the scenario of switching your choice.
\vspace{3cm}
\end{exmp}
\subsection{Probability Theory}
\begin{defn}
We can assign other probabilities to events, the two conditions that must be met are, that $P(S)=1$ and $0\leq P(T)\leq 1$ for all $T\subseteq S$.
\newline
In \textbf{discrete} probability theory, we call the elements of $S$ 'outcomes'. We say that the probability of an event $E\subseteq S$ is the sum of the probabilities of the outcomes. (so we really only assign probabilities to the outcomes, and the probability of the events is a result of the outcome probabilities). When all outcomes have the same probability we call this the \textbf{uniform} distribution. Here the probabilities are uniformly distributed among the outcomes.
\end{defn}
\begin{exmp}
We can model the rolling of a fair die, by defining $S=\{1,2,3,4,5,6\}$ and assigning a uniform distribution. 
\newline
Compute the probabilities of the following events. 
\newline
\begin{itemize}
\item We roll a 1
\vspace{1cm}
\item We roll an even
\vspace{1cm}
\end{itemize}
\end{exmp}
\begin{defn}{Conditional Probability}
Let $E$ and $F$ be events with $p(F) > 0$. The conditional probability of $E$ given $F$, denoted by $p(E | F)$, is defined as
\[
p(E | F) = \frac{p(E \cap F)}{p(F)}
\]
\end{defn}
\begin{exmp}What is the probability that we rolled a 3 if we know we rolled an odd?
\vspace{1cm}
\end{exmp}
\begin{defn}[Independence]
The events $E$ and $F$ are \textbf{independent} if and only if $P(E \cap F)=P(E)\cdot P(F)$
\end{defn}
\begin{exmp}Prove that if $E$ and $F$ are independent, then
\newline
$P(E|F)=P(E)$ and $P(F|E)=P(F)$
\vspace{2cm}

\end{exmp}
\begin{exmp} Suppose that we flip $n$ fair coins and keep track of whether they land on heads of tails. 
Let $H_1$ be the event that the first coin lands on heads, $T_1$ be the event the first coin lands on tails. More generally, let $H_n$ be the event that the $n$th coin lands on heads, $T_n$ be the event the $n$th coin lands on tails and 
\newline
Find the following probabilities:
\begin{itemize}
\item $P(H_1)$
\vspace{1cm}
\item $P(T_4)$
\vspace{1cm}
\item $P(H_1\cap H_2)$
\vspace{1cm}
\item $P(H_1)\cdot P(H_2)$
\vspace{1cm}
\item $P(H_1|H_2)$
\vspace{1cm}
\end{itemize}


\end{exmp}
\begin{defn}[Bernoulli trials]
Consider an experiment with only two possible outcomes, such as generating a random bit (0 or 1) or flipping a coin (heads or tails). Each iteration of such an experiment, where there are only two potential results, is termed a Bernoulli trial, (named after James Bernoulli).

A possible outcome of a Bernoulli trial is classified as either a \textbf{success} or a \textbf{failure}. Let $p$ represent the probability of success and $q$ represent the probability of failure, with the constraint that $p + q = 1$.

For many problems, one can find solutions by determining the probability of achieving k successes when an experiment consists of n mutually independent Bernoulli trials. The term "mutually independent" refers to the condition where the probability of success in any given trial remains p, irrespective of any information about the outcomes of the other trials.
\end{defn}
\begin{exmp}[Fair Coin]$ $\\
What is the probability that after 2 coin flips exactly one is heads?
\newline
\vspace{1cm}
\newline
What is the probability that after 4 coin flips exactly 2 are heads?
\newline
\vspace{1cm}
\newline
What is the probability that after 6 coin flips exactly 2 are heads?
\newline
\vspace{1cm}

\end{exmp}

\noindent \textbf{Recall} If $A\cap B=\emptyset$ then $P(A \cup B)=P(A)+P(B)$

\noindent \textbf{Notice} $P(HT)=P(TH)$
\begin{exmp}[Unfair Coin]
Suppose that we have an unfair coin which lands on heads with probability $.7$ and tails with probability $.3$.
\\
What is the probability that after 2 coin flips exactly one is heads?
\newline
\vspace{1cm}
\newline
What is the probability that after 4 coin flips exactly 2 are heads?
\newline
\vspace{1cm}
\newline
What is the probability that after 6 coin flips exactly 2 are heads?
\newline
\vspace{1cm}


\end{exmp}
\begin{defn}
The probability of exactly \(k\) successes in \(n\) independent Bernoulli trials, with a probability of success \(p\) and a probability of failure \(q = 1 - p\), is given by the binomial probability formula:

\[
C(n, k) p^k q^{n-k}
\]

where \(X\) is the random variable representing the number of successes, \(C(n, k)\) is the binomial coefficient, and \(p\) and \(q\) are the probabilities of success and failure, respectively.

\end{defn}

\begin{defn}[Random Variable]
A \textbf{random variable} is a function from the sample space of an experiment to the set of real numbers. That is, a random variable assigns a real number to each possible outcome. (Equivalently, a random variable is a function that maps each event to a real number.)
\end{defn}

\begin{exmp}
Suppose that a coin is flipped three times. Let \(X(t)\) be the random variable that equals the
number of heads that appear when \(t\) is the outcome. Then \(X(t)\) takes on the following values:
\newline
\vspace{2cm}
\end{exmp}

\begin{defn}[Distribution]
The \textbf{distribution} of a random variable \(X\) on a sample space \(S\) is the set of pairs 
\[
(r,\, P(X = r))
\]
for all \(r \in X(S)\), where \(P(X = r)\) is the probability that \(X\) takes the value \(r\). 
The distribution is completely determined by the probabilities \(P(X = r)\) for each \(r \in X(S)\).
\end{defn}



\begin{exmp}[Distribution of a Dice Roll]
Let the sample space be
\[
S = \{1, 2, 3, 4, 5, 6\},
\]
corresponding to the outcomes of rolling a fair six-sided die.  
Define the random variable \(X(t)\) to be the number showing on the upper face.

Then the distribution of \(X\) is
\[
\{(1, \tfrac{1}{6}), (2, \tfrac{1}{6}), (3, \tfrac{1}{6}), (4, \tfrac{1}{6}), (5, \tfrac{1}{6}), (6, \tfrac{1}{6})\}.
\]

That is, for each possible value \(r \in X(S)\),
\[
P(X = r) = \tfrac{1}{6}.
\]

This example illustrates a \emph{discrete uniform distribution}—each outcome in the range of \(X\) is equally likely.
\newline
\vspace{1cm}
\end{exmp}



\subsection{Bayes' Theorem}
\begin{defn}
\textbf{Bayes' Theorem:} Suppose that \(E\) and \(F\) are events from a sample space \(S\) such that \(p(E) \neq 0\) and \(p(F) \neq 0\). Then
\[
p(F | E) = \frac{p(E | F) \cdot p(F)}{p(E | F) \cdot p(F) + p(E | \overline{F}) \cdot p(\overline{F})}
\]
\end{defn}
\begin{exmp}[Bayes use in real life]

Let:
\begin{align*}
& P(\text{Disease}) = 0.01 \quad (\text{1\% of the population has the disease}) \\
& P(\text{Positive} | \text{Disease}) = 0.95 \quad (\text{95\% chance of a positive test if the person has the disease}) \\
& P(\text{Negative} | \text{No Disease}) = 0.90 \quad (\text{90\% chance of a negative test if the person does not have the disease}) \\
& P(\text{No Disease}) = 1 - P(\text{Disease}) = 0.99 \quad (\text{99\% of the population does not have the disease})
\end{align*}

The probability of having the disease given a positive test result using Bayes' Theorem is given by:
\[
P(\text{Disease} | \text{Positive}) = \frac{P(\text{Positive} |\text{Disease}) \cdot P(\text{Disease})}{P(\text{Positive} | \text{Disease}) \cdot P(\text{Disease}) + P(\text{Negative} | \text{No Disease}) \cdot P(\text{No Disease})}
\]

Substituting the given values:
\[
P(\text{Disease} | \text{Positive}) = \frac{0.95 \cdot 0.01}{0.95 \cdot 0.01 + (1 - 0.90) \cdot 0.99}
\]

\end{exmp}
\subsection{Expected Value}
\begin{defn}
The expected value, also called the expectation or mean, of the random variable \(X\) on the sample space \(S\) is equal to
\[ E(X) = \sum_{s \in S} p(s)X(s). \]

The deviation of \(X\) at \(s \in S\) is \(X(s) - E(X)\), the difference between the value of \(X\) and the mean of \(X\).

\end{defn}
\begin{exmp}
Let $X$ be the number that comes up when a fair die is rolled. What
is the expected value of $X$?
\vspace{3cm}
\end{exmp}

\begin{thm}
If \(X_i, i = 1, 2, \ldots, n\) with \(n\) as a positive integer, are random variables on \(S\), and if \(a\) and \(b\) are real numbers, then
\begin{enumerate}
    \item \(E(X_1 + X_2 + \ldots + X_n) = E(X_1) + E(X_2) + \ldots + E(X_n)\)
    \item \(E(aX + b) = aE(X) + b\).
\end{enumerate}

\end{thm}
\begin{exmp}
Find the expected number of successes when $n$ mutually independent Bernoulli trials are performed,
where $p$ is the probability of success on each trial.
\end{exmp}
% \begin{exmp}
% The expected number of successes when n mutually independent Bernoulli trials are performed,
% where p is the probability of success on each trial?
% \vspace
% \end{exmp}




\end{document}
